\documentclass[11pt,a4paper]{article}
\usepackage[margin=2.4cm]{geometry}
\usepackage{amsmath,amssymb,amsthm,mathtools}
\usepackage[protrusion=true,expansion=false]{microtype}
\usepackage{booktabs,enumitem}
\usepackage{placeins}
\usepackage{tikz}
\usetikzlibrary{arrows.meta,positioning,calc,fit,backgrounds}
\definecolor{tierspend}{RGB}{200,60,60}
\definecolor{tierfull}{RGB}{210,140,40}
\definecolor{tierin}{RGB}{60,130,180}
\definecolor{tierrecv}{RGB}{70,150,90}
\usepackage[colorlinks=true,linkcolor=blue!50!black,citecolor=blue!50!black,urlcolor=blue!50!black]{hyperref}

\emergencystretch=3em
\hbadness=10000
\hfuzz=2pt

\newtheorem{theorem}{Theorem}[section]
\newtheorem{lemma}[theorem]{Lemma}
\newtheorem{proposition}[theorem]{Proposition}
\newtheorem{corollary}[theorem]{Corollary}
\theoremstyle{definition}
\newtheorem{definition}[theorem]{Definition}
\newtheorem{construction}[theorem]{Construction}
\newtheorem{assumption}[theorem]{Assumption}
\newtheorem{remark}[theorem]{Remark}
\newtheorem{example}[theorem]{Example}

\newcommand{\F}{\mathbb{F}}
\newcommand{\Z}{\mathbb{Z}}
\newcommand{\N}{\mathbb{N}}
\newcommand{\G}{\mathbb{G}}
\newcommand{\E}{\mathbb{E}}
\newcommand{\PP}{\mathbb{P}}
\renewcommand{\Pr}{\mathrm{Pr}}
\newcommand{\Fp}{\F_p}
\newcommand{\Fq}{\F_q}
\newcommand{\ip}[2]{\langle #1,\, #2 \rangle}
\newcommand{\Comm}{\mathsf{Commit}}
\newcommand{\Open}{\mathsf{Open}}
\newcommand{\Setup}{\mathsf{Setup}}
\newcommand{\Verify}{\mathsf{Verify}}
\newcommand{\negl}{\mathsf{negl}}
\newcommand{\poly}{\mathsf{poly}}
\newcommand{\PPT}{\mathsf{PPT}}
\newcommand{\Adv}{\mathcal{A}}
\newcommand{\Prov}{\mathcal{P}}
\newcommand{\Ver}{\mathcal{V}}
\newcommand{\Sim}{\mathcal{S}}
\newcommand{\Ext}{\mathcal{E}}
\newcommand{\Rel}{\mathcal{R}}
\newcommand{\Lang}{\mathcal{L}}
\newcommand{\nfsym}{\mathsf{nf}}
\newcommand{\cmsym}{\mathsf{cm}}
\newcommand{\cmx}{\mathsf{cmx}}
\newcommand{\cvsym}{\mathsf{cv}}
\newcommand{\rk}{\mathsf{rk}}
\newcommand{\Pallas}{\ensuremath{\mathsf{Pallas}}}
\newcommand{\Vesta}{\ensuremath{\mathsf{Vesta}}}
\newcommand{\Ep}{\ensuremath{\mathcal{E}}}
\newcommand{\NoteCommit}{\ensuremath{\mathsf{NoteCommit}}}
\newcommand{\Sinsemilla}{\ensuremath{\mathsf{Sinsemilla}}}
\newcommand{\ShortCommit}{\ensuremath{\mathsf{ShortCommit}}}
\newcommand{\CommitIvk}{\ensuremath{\mathsf{Commit}^{\mathsf{ivk}}}}
\newcommand{\Extract}{\ensuremath{\mathsf{Extract}}}
\newcommand{\Acc}{\ensuremath{\mathsf{Acc}}}
\newcommand{\GroupHash}{\ensuremath{\mathsf{GroupHash}}}
\newcommand{\ask}{\ensuremath{\mathsf{ask}}}
\newcommand{\aksym}{\ensuremath{\mathsf{ak}}}
\newcommand{\nksym}{\ensuremath{\mathsf{nk}}}
\newcommand{\ivk}{\ensuremath{\mathsf{ivk}}}
\newcommand{\fvk}{\ensuremath{\mathsf{fvk}}}
\newcommand{\ovk}{\ensuremath{\mathsf{ovk}}}
\newcommand{\dk}{\ensuremath{\mathsf{dk}}}
\newcommand{\pkd}{\ensuremath{\mathsf{pk_d}}}
\newcommand{\gd}{\ensuremath{\mathsf{g_d}}}
\newcommand{\rcv}{\ensuremath{\mathsf{rcv}}}
\newcommand{\rcm}{\ensuremath{\mathsf{rcm}}}
\newcommand{\rivk}{\ensuremath{\mathsf{r_{ivk}}}}
\newcommand{\rsk}{\ensuremath{\mathsf{rsk}}}

\renewenvironment{abstract}{%
  \small\begin{center}{\bfseries\abstractname}\end{center}\medskip\noindent\ignorespaces
}{\par\medskip}

\title{\textbf{\Huge Sync Guide}\\[6pt]\large Light-client synchronisation: compact blocks, the service, and the scanning pipeline}
\author{m@rek.onl}
\date{}
\begin{document}
\maketitle
\thispagestyle{empty}
\begin{abstract}
Assuming the Orchard protocol of the \emph{Orchard Guide} and the wallet layer
of the \emph{Wallet Guide}, this volume of \emph{The Zcash Arboretum}
documents how a light client synchronises: the compact-block format (ZIP-307)
and its deployed divergences, the light-client service surface as lightwalletd
and Zaino implement it, the scanning pipeline of \textsf{librustzcash} with
its verify-first reorg discipline, commitment-tree synchronisation, and ---
stated without euphemism --- what the server learns about its clients. The
wire format is deployed but underspecified; where the draft ZIP and the
deployed implementations disagree, this volume documents the deployment and
flags the divergence.
\end{abstract}
\clearpage
\tableofcontents
\clearpage
\section{Compact blocks (ZIP-307)}\label{sec:sync-compact}

A wallet that trusts no server detects its notes by trial-decrypting every
shielded output on the chain (\emph{Orchard Guide}, \S``Transmission and
detection of a note: encryption, trial decryption, and synchronisation'').
That obligation costs twice: one key agreement of computation per output, and
the bandwidth of downloading every output. ZIP~307 (\emph{Light Client
Protocol for Payment Detection}) attacks the bandwidth half. A \emph{compact
block} is ``a packaging of ONLY the data from a block needed to'' (1)~detect
a payment to one's shielded address, (2)~detect a spend of one's notes, and
(3)~update witnesses for new spend proofs (ZIP~307, ``Compact Stream
Format''). The deployed protocol widens this beyond ZIP~307's Sapling-only
scope to every shielded pool, and adds a fourth purpose: detecting spends of
the coins of the wallet's transparent addresses
(\texttt{compact\_formats.proto}, header comment, lines~21--26). The
computational half survives compaction --- every compact output must still be
trial-decrypted --- and that residual linear scan is among the costs whose
removal motivates the successor design (\emph{Tachyon Guide}, \S``What the
fold removes'').

%% TODO: the service API surface (service.proto, CompactTxStreamer:
%% GetBlockRange, GetTreeState, GetTransaction, GetSubtreeRoots) is deferred
%% to its own section; add \S\ref once that section exists. ZIP-307's own
%% "Proxy operation" sketch (RangeFilter, FullTransaction, ChainSpec) was
%% never deployed under those names and the ZIP warns its API "may differ
%% from the implementation" -- cover the divergence there, not here.

\subsection{Normative status and wire-format history}%
\label{sec:sync-compact-status}

ZIP~307 carries Status: Draft, covers Sapling only, and nowhere mentions
Orchard, \texttt{ChainMetadata}, or transparent compact data
(\texttt{zip-0307.rst}, header and whole file). Its message definitions are
stale relative to deployment: the ZIP text defines
\texttt{CompactBlock\{BlockID id = 1; vtx = 3\}} and
\texttt{CompactTx\{txIndex, txHash, spends, outputs\}}, with element messages
named \texttt{CompactSpend} and \texttt{CompactOutput} (ZIP~307, ``Compact
Stream Format''), whereas the deployed messages of
\S\ref{sec:sync-compact-wire} use different names, field numbers, and fields.
The renames to \texttt{CompactSaplingSpend}/\texttt{CompactSaplingOutput}
date to \textsf{lightwallet-protocol} v0.3.1 (2021-12-09), which also added
\texttt{CompactOrchardAction} (\textsf{lightwallet-protocol},
\texttt{CHANGELOG.md}).

The canonical proto files live in the \texttt{zcash/lightwallet-protocol}
repository and are vendored --- git subtree or symlink --- into
\textsf{lightwalletd}, \textsf{librustzcash} (the files under
\texttt{zcash\_client\_backend/proto/} are symlinks into
\texttt{lightwallet-protocol/walletrpc/}), and \textsf{Zaino}. The
repository's README states that it ``IS not a specification'' and defers to
ZIP~307 (\textsf{lightwalletd},
\texttt{lightwallet-protocol/README.md}) --- which is itself stale. No
current document therefore normatively specifies the deployed wire format.
Per our conventions we classify the compact-block format \emph{deployed but
underspecified}: we cite \textsf{lightwallet-protocol} tagged releases as the
de facto definition, and flag each point where ZIP text and deployment
disagree (\S\ref{sec:sync-compact-divergences}).

\begin{table}[htbp]
\centering\small
\begin{tabular}{@{}llp{8.2cm}@{}}
\toprule
Tag & Date & Change \\
\midrule
v0.3.1 & 2021-12-09 & \texttt{CompactOrchardAction} added;
  \texttt{CompactSpend}/\texttt{CompactOutput} renamed
  \texttt{CompactSaplingSpend}/\texttt{CompactSaplingOutput} \\
v0.3.2 & 2021-12-09 & \texttt{CompactOrchardAction.encCiphertext} renamed
  \texttt{ciphertext} \\
v0.3.5 & 2023-07-03 & \texttt{ChainMetadata} and
  \texttt{CompactBlock.chainMetadata} added; \texttt{GetSubtreeRoots}/%
  \texttt{SubtreeRoot} added; nullifier RPCs added;
  \texttt{CompactSaplingOutput.epk} renamed \texttt{ephemeralKey} \\
v0.4.0 & 2025-12-03 & \texttt{CompactTxIn}/\texttt{TxOut} and
  \texttt{vin}/\texttt{vout} added; \texttt{BlockRange.poolTypes} added;
  \texttt{LightdInfo.lightwalletProtocolVersion} added;
  \texttt{CompactTx.hash} renamed \texttt{txid}
  (serialization-compatible) \\
v0.4.1 & 2026-02-20 & nullifier RPCs designated deprecated in the
  proto \\
v0.5.0 & 2026-06-30 & Ironwood fields added; \texttt{protoVersion} removed,
  field number~1 and name reserved \\
\bottomrule
\end{tabular}
\caption{Wire-format history of \textsf{lightwallet-protocol}
(\texttt{CHANGELOG.md}, librustzcash-vendored copy).}%
\label{tab:sync-lwp-history}
\end{table}

The three deployed codebases vendor three different revisions:
\textsf{lightwalletd} v0.4.1 (\texttt{protoVersion} present, no Ironwood);
\textsf{librustzcash} main v0.5.0 (\texttt{protoVersion} removed and
reserved; Ironwood fields present); \textsf{Zaino} an intermediate snapshot
carrying \emph{both} the Ironwood fields and \texttt{protoVersion} (diff of
the three \texttt{compact\_formats.proto} copies). We flag the skew, and a
changelog-date contradiction, in \S\ref{sec:sync-compact-divergences}.
Deployed-behaviour claims in this section are verified against
\textsf{librustzcash} \texttt{e30517e4}, \textsf{lightwalletd}
\texttt{61fee32e}, \textsf{Zaino} \texttt{0e057e22}, and \textsf{zips}
\texttt{6961098} (checkouts of 2026-07-14).

\subsection{The deployed messages}\label{sec:sync-compact-wire}

The messages live in the protobuf package
\texttt{cash.z.wallet.sdk.rpc} and declare \texttt{proto3} syntax
(\texttt{compact\_formats.proto}, lines~5--11).
Proto3 semantics make every field optional: an absent field
decodes to zero or empty, indistinguishable from an explicit zero. The
scanner must defend against this --- the default-zero tree sizes of
\S\ref{sec:sync-compact-scan} are the load-bearing case.

\begin{center}\footnotesize
\begin{tabular}{@{}p{8.8cm}p{4.6cm}@{}}
\multicolumn{2}{@{}l}{\ttfamily message CompactBlock \{} \\
\quad\ttfamily uint32 protoVersion = 1; & removed in v0.5.0; number and
  name reserved \\
\quad\ttfamily uint64 height = 2; & block height \\
\quad\ttfamily bytes hash = 3; & block hash, 32 bytes \\
\quad\ttfamily bytes prevHash = 4; & parent block hash, 32 bytes \\
\quad\ttfamily uint32 time = 5; & Unix epoch time \\
\quad\ttfamily bytes header = 6; & ``should always be unset (empty)'' \\
\quad\ttfamily repeated CompactTx vtx = 7; & compact transactions \\
\quad\ttfamily ChainMetadata chainMetadata = 8; & end-of-block tree sizes \\
\multicolumn{2}{@{}l}{\ttfamily \}} \\[4pt]
\multicolumn{2}{@{}l}{\ttfamily message ChainMetadata \{} \\
\quad\ttfamily uint32 saplingCommitmentTreeSize = 1; & Sapling tree size at
  end of block \\
\quad\ttfamily uint32 orchardCommitmentTreeSize = 2; & Orchard tree size at
  end of block \\
\quad\ttfamily uint32 ironwoodCommitmentTreeSize = 3; & v0.5.0 only \\
\multicolumn{2}{@{}l}{\ttfamily \}} \\[4pt]
\multicolumn{2}{@{}l}{\ttfamily message CompactTx \{} \\
\quad\ttfamily uint64 index = 1; & position within the block \\
\quad\ttfamily bytes txid = 2; & 32 bytes, protocol byte order \\
\quad\ttfamily uint32 fee = 3; & unpopulated in practice \\
\quad\ttfamily repeated CompactSaplingSpend spends = 4; & \\
\quad\ttfamily repeated CompactSaplingOutput outputs = 5; & \\
\quad\ttfamily repeated CompactOrchardAction actions = 6; & \\
\quad\ttfamily repeated CompactTxIn vin = 7; & coinbase input omitted \\
\quad\ttfamily repeated TxOut vout = 8; & \\
\quad\ttfamily repeated CompactOrchardAction ironwoodActions = 9; &
  v0.5.0 only \\
\multicolumn{2}{@{}l}{\ttfamily \}} \\
\end{tabular}
\end{center}

The container messages, from the deployed protos (\textsf{lightwalletd}
vendored copy, \texttt{walletrpc/}\allowbreak%
\texttt{compact\_formats.proto}, lines~14--80; \textsf{librustzcash} copy,
lines~14--86):

\begin{itemize}
\item Field \texttt{header} ``should always be unset (empty)'' per the
  proto comment (lines~27--30). ZIP~307's block
  header validation --- Equihash, \texttt{bits}, difficulty --- is therefore
  unimplementable against deployed servers; the ZIP itself flags the section
  as ``only partially implemented: currently only \texttt{prevHash} is
  checked'' (ZIP~307, ``Block header validation''). Client-side, only
  \texttt{prevHash} and height continuity are enforced
  (\S\ref{sec:sync-compact-scan}).
%% TODO: confirm no deployed client ever requests headers before stating the
%% header path dead categorically (open question).
\item Field \texttt{protoVersion} has three deployed behaviours:
  \textsf{lightwalletd} never sets it, the assignment commented out as a
  \texttt{TODO} (\texttt{parser/block.go}, line~112); \textsf{Zaino} sets
  it to~1 (\texttt{zaino-fetch}, \texttt{src/chain/block.rs}, line~218);
  v0.5.0 removed it, reserving field number~1 and the name
  (\textsf{librustzcash} copy, lines~33--36). The \textsf{librustzcash}
  scanner never reads it. Flagged in
  \S\ref{sec:sync-compact-divergences}.
\item Message \texttt{ChainMetadata} states each pool's note commitment tree
  size \emph{as of the end of this block}, as \texttt{uint32}
  (\texttt{compact\_formats.proto}, lines~14--17): the tree holds at most
  $2^{32}$ leaves, and the scanner maps arithmetic overflow to
  \texttt{ScanError::TreeSizeOverflow} (\textsf{zcash\_client\_backend},
  \texttt{src/scanning.rs}, lines~649--654). Section
  \ref{sec:sync-compact-scan} shows how these sizes turn a compact block
  into note positions.
\item Field \texttt{txid} carries the 32-byte transaction identifier in
  protocol byte order; the proto comment mandates that it ``MUST NOT be
  reversed or hex-encoded'' --- the reversed hex form is display-only ---
  citing protocol spec \S7.1.1, ``Transaction Identifiers''
  (\texttt{compact\_formats.proto}, lines~52--57).
\item The \texttt{fee} field is documented ``present if server can
  provide'', but neither deployed server populates it:
  \textsf{lightwalletd} leaves it commented out with \texttt{TODO:
  calculate fees} (\texttt{parser/transaction.go}, line~407) and
  \textsf{Zaino} hardcodes \texttt{fee: 0} (\texttt{zaino-fetch},
  \texttt{src/chain/transaction.rs}, line~369).
  Clients cannot rely on it. For a transaction with no transparent inputs
  the proto comment gives the client-side reconstruction
  (\texttt{compact\_formats.proto}, lines~59--64):
  \[
  \mathrm{fee} \;=\; \mathrm{valueBalanceSapling}
    + \mathrm{valueBalanceOrchard}
    + \textstyle\sum \mathrm{vPubNew}
    - \sum \mathrm{vPubOld}
    - \sum \mathrm{tOut},
  \]
  with $\mathrm{valueBalanceIronwood}$ joining the sum in
  v0.5.0.\footnote{Type \texttt{uint32} additionally caps the representable
  fee at $2^{32}-1$ zatoshi ($\approx 42.9$~ZEC); the proto comment does not
  remark on this.}
\item Field \texttt{vin} omits the coinbase transaction's single
  null-outpoint input; a light client recognises a coinbase
  \texttt{CompactTx} by $\texttt{index} = 0$, coinbase being always first in
  a block. Server \textsf{lightwalletd} implements exactly this, skipping
  \texttt{vin} when $\texttt{index} = 0$ (\texttt{compact\_formats.proto},
  lines~70--76; \texttt{parser/transaction.go}, lines~398--403). The
  transparent element messages are
  $\texttt{CompactTxIn}\{\texttt{prevoutTxid},\ \texttt{prevoutIndex}\}$ and
  $\texttt{TxOut}\{\texttt{value},\ \texttt{scriptPubKey}\}$
  (\texttt{compact\_formats.proto}, lines~82--104).
\end{itemize}

\paragraph{The shielded elements.} Each shielded element message keeps only
the detection-relevant bytes of its full on-chain description.

\begin{table}[htbp]
\centering\footnotesize
\begin{tabular}{@{}lp{5.4cm}rrr@{}}
\toprule
Message & Fields (bytes) & Payload & Full form & Saving \\
\midrule
\texttt{CompactSaplingSpend} & \texttt{nf} (32) & 32 & 384 &
  $\sim$90\% \\
\texttt{CompactSaplingOutput} & \texttt{cmu} (32),
  \texttt{ephemeralKey} (32), \texttt{ciphertext} (52) & 116 & $580+32$ &
  $\sim$80\% \\
\texttt{CompactOrchardAction} & \texttt{nullifier} (32), \texttt{cmx} (32),
  \texttt{ephemeralKey} (32), \texttt{ciphertext} (52) & 148 & --- & --- \\
\bottomrule
\end{tabular}
\caption{Compact shielded elements (\texttt{compact\_formats.proto},
lines~106--130). The Sapling payload and saving figures are ZIP~307's
(``Spend Compression'', ``Output Compression''); the 148-byte Orchard total
is arithmetic, not stated in the sources. Payload figures count raw bytes;
on the wire, a \texttt{CompactSaplingSpend} encodes to 34 bytes (36 as a
\texttt{CompactTx.spends} element), a \texttt{CompactSaplingOutput} to 122
(124), and a \texttt{CompactOrchardAction} to 156 (159 --- its 156-byte
payload needs a two-byte length varint). Field headers add 2 bytes per
32-byte field and 2 per 52-byte ciphertext. Computed by script,
proto-encoding every compact element of mainnet block $3{,}412{,}934$ (3
spends, 8 outputs, 10 actions; \texttt{GetBlock}, lightwalletd v0.4.19,
2026-07-15); the sizes are block-independent because all field lengths are
fixed.}\label{tab:sync-compact-sizes}
\end{table}

A \texttt{CompactSaplingSpend} retains only the 32-byte nullifier of the
384-byte Spend description ($\cvsym$~32, anchor~32, nullifier~32, $\rk$~32,
zkproof~192, spendAuthSig~64) --- the one field needed to detect a spend of
one's own notes; ZIP~307 quotes this as a 90\% bandwidth improvement
(ZIP~307, ``Spend Compression''). A \texttt{CompactSaplingOutput} retains
$\mathsf{cmu}$, the ephemeral key, and the first 52 bytes of
$\mathsf{encCiphertext}$, ``Total size is 116 bytes'' per the proto comment,
against 580 bytes of ciphertext plus the 32-byte ephemeral key streamed in
full form --- ZIP~307's 80\% reduction (ZIP~307, ``Output Compression'';
\textsf{lightwalletd} \texttt{compact\_formats.proto}, lines~114--122; full
encodings per protocol spec \S\S7.3--7.4). A \texttt{CompactOrchardAction}
serves both roles of an Action in one message: its \texttt{nullifier}
detects spends of the input note, while \texttt{cmx},
\texttt{ephemeralKey}, and \texttt{ciphertext} detect and place the output
note (\texttt{compact\_formats.proto}, lines~124--130; full Action encoding
per protocol spec \S7.5).

\subsection{Omissions and fetch-on-detect}\label{sec:sync-compact-omitted}

Compact blocks omit the 512-byte memo, the 16-byte AEAD tag, the 80-byte
$\mathsf{outCiphertext}$, $\cvsym$, $\rk$, anchors, zkproofs, and signatures
(ZIP~307, ``Transaction privacy''; sizes per
\textsf{zcash\_note\_encryption}~0.4.1, \texttt{src/lib.rs}, lines~44--57).
Recovering a memo, or spend information via the outgoing ciphertext, thus
requires downloading the full transaction --- deployed as the
\texttt{GetTransaction} RPC. ZIP~307 attaches three requirements to that
fetch: the client SHOULD obscure its interest by also downloading numerous
uninteresting transactions, SHOULD download all transactions of any block it
touches, and MUST convey the privacy reduction to the user (ZIP~307,
``Transaction privacy'').

The deployed pipeline expresses the fetch as data: scanning emits
\texttt{TransactionDataRequest::}\allowbreak\texttt{Enhancement(txid)} ---
``download of complete raw transaction data'' --- which the caller
satisfies via \texttt{GetTransaction} and feeds to
\texttt{wallet::}\allowbreak\texttt{decrypt\_and\_store\_transaction}
(\textsf{zcash\_client\_backend}, \texttt{src/data\_api.rs},
lines~1185--1207, 2210). The crate's own sync loop documents that it does
not yet perform enhancement (\textsf{zcash\_client\_backend},
\texttt{src/sync.rs}, lines~1--10).

\subsection{Compact trial decryption}\label{sec:sync-compact-decrypt}

The 52-byte \texttt{ciphertext} field is the prefix of
$\mathsf{encCiphertext}$ that encrypts exactly the \emph{compact note
plaintext}: lead byte (1)~$\|$ diversifier $d$~(11) $\|$ value $v$~(8) $\|$
$\mathsf{rseed}$~(32), so $\mathrm{COMPACT\_NOTE\_SIZE} = 52$; the full
plaintext appends the 512-byte memo, and the full ciphertext a 16-byte tag,
giving $580$ bytes (\textsf{zcash\_note\_encryption}~0.4.1,
\texttt{src/lib.rs}, lines~44--57). ZIP~307's rationale: these bytes
``contain the contents and opening of the note commitment, which is all of
the data needed to spend the note and to verify that the note is spendable''
(ZIP~307, ``Output Compression''). The \emph{Orchard Guide} constructs
compact decryption for Orchard in \S``Transmission and detection of a note:
encryption, trial decryption, and synchronisation''; we do not repeat the
construction, and record here only the deployed mechanics and the integrity
argument.

Function \texttt{try\_compact\_note\_decryption} (doc-comment: ``Implements
the procedure specified in ZIP 307'') proceeds as follows
(\textsf{zcash\_note\_encryption}~0.4.1, \texttt{src/lib.rs},
lines~567--604):
\begin{align*}
K &:= \mathsf{KDF}\bigl(\mathsf{KA.Agree}(\ivk, \mathsf{epk}),\,
      \mathsf{ephemeralKey}\bigr),\\
\mathsf{np} &:= \mathsf{ciphertext}[0..52] \oplus
      \mathsf{ChaCha20}_{K}\bigl(\mathrm{nonce} = 0^{96}\bigr)[64..116],
\end{align*}
a bare ChaCha20 keystream with zero nonce, sought to byte offset~$64$ ---
skipping keystream block~0, which the full AEAD reserves for the Poly1305
key (\texttt{src/lib.rs}, lines~590--593). The candidate $\mathsf{np}$ is
parsed as a compact note plaintext (no memo), then
\texttt{check\_note\_validity} accepts iff
\begin{enumerate}[label=(\roman*)]
\item the plaintext parses, and
\item the recomputed commitment matches the on-chain one,
  $\Extract\bigl(\NoteCommit_{\rcm}(\ldots)\bigr) = \mathsf{cmu}$
  (respectively $\cmx$), and
\item per ZIP-212, the ephemeral secret $\mathsf{esk}$ re-derived from
  $\mathsf{rseed}$ reproduces the published key: the byte encoding of
  $\mathsf{KA.DerivePublic}(\mathsf{esk}, \gd)$ equals
  \texttt{ephemeralKey}, compared in constant time
\end{enumerate}
(\texttt{src/lib.rs}, lines~524--604). Condition~(ii) is the integrity
substitute: truncation at 52 bytes discards the AEAD tag, so a compact
decryption cannot be authenticated via Poly1305; ZIP~307 instead recomputes
the note commitment from the decrypted plaintext and compares against the
on-chain $\mathsf{cmu}$, itself authenticated by the binding signature over
the transaction (ZIP~307, ``Output Compression''). Condition~(iii) is the
deployed strengthening beyond the tag: it defeats a mauled ephemeral key
(\emph{Orchard Guide}, \S``Transmission and detection of a note'').

\paragraph{Lead byte and the ZIP-212 grace period.} The lead byte is
$\mathtt{0x01}$ pre-Canopy and $\mathtt{0x02}$ under ZIP-212; deployed code
switches on \texttt{zip212\_enforcement}: during the grace period of
$32{,}256$ blocks after Canopy activation both lead bytes are accepted, and
thereafter only $\mathtt{0x02}$ (\textsf{zcash\_primitives},
\texttt{src/transaction/components/sapling.rs}, lines~27--40;
\textsf{zcash\_protocol}, \texttt{src/consensus.rs}, line~681). ZIP~307's
pseudocode disagrees with itself here: its two ``[Canopy onward]'' lines are
\emph{both} conditioned on $\mathsf{height} < \mathsf{CanopyActivationHeight}
+ \mathsf{ZIP212GracePeriod}$, the first rejecting $\mathsf{leadByte} \notin
\{\mathtt{0x01}, \mathtt{0x02}\}$ and the second rejecting
$\mathsf{leadByte} \neq \mathtt{0x02}$ (ZIP~307, ``Scanning for relevant
transactions'', \texttt{zip-0307.rst}, lines~275--276). Read literally, the
pair rejects $\mathtt{0x01}$ \emph{during} the grace period and constrains
nothing after it --- the inverse of the intent; the second condition should
read~$\geq$. Deployed code implements the intended semantics. The defect
is tracked in \texttt{research/upstream-defects.md}, item~1 (ZIP~307:
second Canopy-onward lead-byte condition rejects $\mathtt{0x01}$ during
the grace period; verified 2026-07-15 at \textsf{zips}
\texttt{69610984}). We flag the
divergence in \S\ref{sec:sync-compact-divergences}.

\subsection{Scanning: spends, positions, and tree sizes}%
\label{sec:sync-compact-scan}

\paragraph{Spend detection.} The scanner parses every compact nullifier
field up front: each \texttt{CompactSaplingSpend.nf} and each
\texttt{CompactOrchardAction.}\allowbreak\texttt{nullifier}. A malformed
length yields \texttt{ScanError::}\allowbreak\texttt{EncodingInvalid}
rather than a panic. The parsed nullifiers are compared against the
wallet's known-unspent nullifiers in constant time (\texttt{find\_spent});
unmatched nullifiers are retained in per-block nullifier maps
(\textsf{zcash\_client\_backend}, \texttt{src/scanning/compact.rs},
lines~312--390; \texttt{src/scanning.rs}, lines~826--828).

\paragraph{Tree building and positions.} Every $\mathsf{cmu}$ and $\cmx$ in
the block --- not only the wallet's --- is appended to the wallet's note
commitment tree with \texttt{Retention} marks, and each detected note is
marked at its position (\textsf{zcash\_client\_backend},
\texttt{src/scanning.rs}, lines~786--824). The position is computed
arithmetically, never by replaying the chain:
\[
\mathrm{pos}(\text{output } i \text{ of } tx)
  \;=\; \mathrm{start\_tree\_size} \;+\; \mathrm{offset}(tx) \;+\; i,
\]
where $\mathrm{offset}(tx)$ counts the outputs of earlier transactions in
the block (\texttt{PositionTracker::}\allowbreak%
\texttt{\{sapling,orchard\}\_note\_position}). The
shardtree that consumes these marks is constructed in the \emph{Orchard
Guide}, \S``Proof of ownership of \emph{some} valid note: the commitment
tree''; the observation that a light wallet reads positions directly off the
block appears in \S``Transmission and detection of a note''.

Positions are load-bearing, not merely convenient: the
\texttt{TreeSizeUnknown} doc-comment records that without the tree size it
is ``impossible to construct the nullifier for a detected note'' --- the
Sapling nullifier depends on the note's position, via $\rho =
\mathsf{MixingPedersenHash}(\cmsym, \mathsf{NotePosition})$; Orchard's does
not, but both pools need positions for witnesses
(\textsf{zcash\_client\_backend}, \texttt{src/scanning.rs}, lines~633--639;
protocol spec, ``Computing $\rho$ values and Nullifiers'' and ``Dummy
Notes'').

\paragraph{Where the tree size comes from.} The start size for a block is
resolved by a ladder (\texttt{tree\_sizes\_around},
\textsf{zcash\_client\_backend}, \texttt{src/scanning/compact.rs},
lines~589--696):
\[
\mathrm{start}(\mathrm{pool}) :=
\begin{cases}
s_{\mathrm{prior}} & \text{prior \texttt{BlockMetadata} known;}\\[2pt]
s_{\mathrm{final}} - n_{\mathrm{block}} &
  \text{\texttt{chainMetadata} present, by checked subtraction;}\\[2pt]
0 & \text{block below the pool's activation height;}\\[2pt]
\bot & \text{otherwise (\texttt{ScanError::TreeSizeUnknown}),}
\end{cases}
\]
where $s_{\mathrm{final}}$ is the \texttt{chainMetadata} end-of-block size
and $n_{\mathrm{block}}$ the block's own output count. The subtraction is
\texttt{checked\_sub}: underflow yields
\texttt{ScanError::TreeSizeInvalid}, which is precisely the guard against
proto3's default-zero metadata (\S\ref{sec:sync-compact-wire}). After
walking all \texttt{CompactTx}s the accumulated position must equal the
\texttt{chainMetadata} size, else scanning fails with
\texttt{ScanError::TreeSizeMismatch\{given, computed\}}
(\texttt{src/scanning/compact.rs}, lines~257--287, 745--770).

The error taxonomy routes recovery: \texttt{TreeSizeMismatch},
\texttt{PrevHashMismatch} (\texttt{prevHash} against the prior block's
hash), and \texttt{BlockHeightDiscontinuity} ($\mathrm{height} \neq
\mathrm{prev}+1$) report $\texttt{is\_continuity\_error()} = \mathrm{true}$
and trigger the truncate-and-rescan reorganisation path (\emph{Wallet
Guide}, \S``Reorganisations: truncate and rescan'');
\texttt{EncodingInvalid}, \texttt{TreeSizeUnknown},
\texttt{TreeSizeInvalid}, and \texttt{TreeSizeOverflow} do not
(\textsf{zcash\_client\_backend}, \texttt{src/scanning.rs},
lines~602--685). The final tree size also drives checkpointing: the helper
\texttt{compact\_tx\_contains\_last\_}\allowbreak%
\texttt{\{sapling,orchard\}\_}\allowbreak\texttt{outputs\_in\_block}
compares the running position plus the current transaction's
output count against the final size, to decide when to write the
block-boundary checkpoint into the commitment tree
(\texttt{src/scanning/compact.rs}, lines~712--731; call sites 404, 444,
482).

\paragraph{The server side of \texttt{ChainMetadata}.}
Server \textsf{lightwalletd} fills the tree sizes from \textsf{zcashd}'s
\texttt{getblock} verbosity-1 response fields \texttt{trees.sapling.size}
and \texttt{trees.orchard.size} (\textsf{lightwalletd},
\texttt{common/common.go}, lines~163--175, 399--400). Upstream,
\textsf{zcashd} computes them by anchor lookup
(\texttt{GetSaplingAnchorAt}/\allowbreak\texttt{GetOrchardAnchorAt} on
\texttt{hashFinalSaplingRoot}/\allowbreak\texttt{hashFinalOrchardRoot})
and omits the
subfield when the anchor is unavailable --- in which case Go's JSON
unmarshal leaves $\texttt{Size} = 0$, the default-zero case
\texttt{TreeSizeInvalid} exists to catch (\textsf{zcashd},
\texttt{src/rpc/blockchain.cpp}, lines~279--297; the \texttt{trees} field
dates to \textsf{zcashd} commit \texttt{cf2b6c796}, ``rpc: Add trees field
to getblock RPC output'').

\subsection{The scan pipeline: chain state, batching, caching}%
\label{sec:sync-compact-pipeline}

Function \texttt{scan\_cached\_blocks} requires a \texttt{ChainState} ---
the final Sapling and Orchard tree frontiers as of
$\mathrm{from\_height} - 1$ --- and asserts $\mathrm{from\_height} =
\texttt{chain\_state.height} + 1$; both trees have depth~32
(\textsf{zcash\_client\_backend}, \texttt{src/data\_api/chain.rs},
lines~504--698). The \texttt{ChainState} comes from the server's
\texttt{GetTreeState} RPC via \texttt{TreeState::to\_chain\_state()}, which
hex-decodes and byte-reverses the block hash and parses empty tree fields as
empty trees (\texttt{src/proto.rs}, lines~367--463). Trial decryption of
compact outputs is batched through \texttt{BatchRunners} with batch
size~100 (\texttt{src/data\_api/chain.rs}, line~642).

The proto accessors split their failure modes: block-level convenience
accessors panic on malformed server data (\texttt{CompactBlock::hash} and
\texttt{prev\_hash} if not 32 bytes, \texttt{height} if outside
\texttt{u32}), while per-output accessors ($\mathsf{cmu}$, $\cmx$,
$\nfsym$, ephemeral key) return
\texttt{CompactFormatError::}\allowbreak%
\texttt{\{InvalidLength, InvalidValue\}}; the
\texttt{From} impls that \emph{build} compact messages from full bundles
take $\mathsf{encCiphertext}[0..\mathrm{COMPACT\_NOTE\_SIZE}]$
(\textsf{zcash\_client\_backend}, \texttt{src/proto.rs}, lines~63--286).
The scanner bounds \texttt{CompactTx.index} to \texttt{u16} (the
\texttt{expect} reads ``Cannot fit more than $2^{16}$ transactions in a
block''), and \texttt{CompactBlock.time} --- \texttt{uint32} Unix epoch,
``if non-zero'' --- flows into \texttt{ScannedBlock} and thence wallet
metadata (\texttt{src/scanning/compact.rs}, lines~314--315, 552--555).

Compact blocks are cached as serialized protos. Crate
\textsf{zcash\_client\_sqlite} offers \texttt{BlockDb}, a
\texttt{compactblocks} SQLite table decoded on read
(\texttt{CompactBlock::decode}), and \texttt{FsBlockDb}, one file per block
plus a \texttt{BlockMeta} row \{height, blockhash, time,
sapling\_outputs\_count, orchard\_actions\_count\} that enables range
queries without decoding (\textsf{zcash\_client\_sqlite},
\texttt{src/chain.rs}, lines~30--230).

ZIP~307's own client-operation sketch --- append every $\mathsf{cmu}$ to a
depth-32 incremental Merkle tree, keep a per-note incremental witness, and
cache roughly 100 recent tree states because ``no full Zcash node will roll
back the chain by more than 100 blocks'' (ZIP~307, ``Creating and updating
note witnesses'') --- describes the pre-shardtree design. The deployed
replacement is the position-marked shardtree fed by \texttt{GetSubtreeRoots}
(RPC added in \textsf{lightwallet-protocol} v0.3.5;
\textsf{zcash\_client\_backend}, \texttt{src/sync.rs}, lines~80--86), whose
mechanics the \emph{Orchard Guide} constructs in \S``Proof of ownership of
\emph{some} valid note: the commitment tree''.

\subsection{Server inclusion semantics and successor surface}%
\label{sec:sync-compact-servers}

Which transactions a compact block contains is itself divergent. ZIP~307
states that ``By default, CompactBlocks only contain CompactTxs for
transactions that contain Sapling spends or outputs'' (ZIP~307, ``Block
header validation''). Server \textsf{lightwalletd} builds and caches
compact representations of \emph{all} transactions, transparent-only ones
included, with \texttt{vin}/\texttt{vout} populated
(\texttt{parser/block.go}, lines~120--125), and \texttt{GetBlock} serves
them unfiltered per the v0.4.1 contract; \texttt{GetBlockRange} filters
per request --- with empty \texttt{poolTypes} it prunes each transaction
to its Sapling and Orchard components and drops transactions left empty
(\texttt{common/common.go},
\texttt{FilterTxPool}/\texttt{filterBlockPool}, lines~526--653) --- so
the default sync stream carries shielded-relevant transactions only, as
Zaino's does. The successor \textsf{Zaino}
filters out \texttt{CompactTx}s whose
compact fields are all empty, and supports per-request pool filtering via
\texttt{PoolTypeFilter} and \texttt{BlockRange.poolTypes}
(\texttt{zaino-fetch/src/chain/block.rs}, lines~193--215) --- a filter added
in v0.4.0, where an empty \texttt{poolTypes} means Sapling and Orchard for
backward compatibility (\texttt{service.proto}, lines~38--41). Both
servers therefore default to shielded-relevant transactions on the sync
stream; we flag the deployed defaults against the ZIP in
\S\ref{sec:sync-compact-divergences}.

Two \textsf{lightwalletd} RPCs, \texttt{GetBlockNullifiers} and
\texttt{GetBlockRangeNullifiers}, strip actions to nullifier-only, drop
outputs, \texttt{vin}, and \texttt{vout}, and zero both
\texttt{ChainMetadata} tree sizes (``these are not needed (we prefer to
save bandwidth)'') (\texttt{frontend/service.go}, lines~212--302). Blocks
so obtained therefore cannot seed the position ladder of
\S\ref{sec:sync-compact-scan} --- zeroed sizes land in the
\texttt{TreeSizeInvalid}/\texttt{TreeSizeUnknown} paths --- so the RPCs
serve spend-status refresh, not scanning. Both were declared deprecated
in the \textsf{lightwallet-protocol} v0.4.0 changelog (2025-12-03);
v0.4.1 designated them deprecated in the proto itself
(\texttt{option deprecated = true}, \texttt{service.proto:247,267}), in
favour of \texttt{GetBlockRange} with \texttt{poolTypes}
(\texttt{service.proto}, lines~244--266;
\texttt{CHANGELOG.md}, v0.4.0/v0.4.1).
%% TODO: whether any deployed client ever consumed the nullifier RPCs for
%% scanning is unverified (open question); the mechanical consequence above
%% follows from the zeroed sizes plus the ladder, not from observed clients.

\paragraph{Ironwood.} The v0.5.0 fields
(\texttt{ironwoodCommitmentTreeSize} $= 3$, \texttt{ironwoodActions} $= 9$)
exist only in \textsf{lightwallet-protocol} v0.5.0;
\textsf{librustzcash} gates the corresponding tree accounting on
\texttt{NetworkUpgrade::Nu6\_3}, and \textsf{lightwalletd} --- the deployed
service --- neither parses nor serves Ironwood data
(\textsf{zcash\_client\_backend}, \texttt{src/scanning/compact.rs},
lines~687--696; \texttt{proto/compact\_formats.proto}, lines~17, 74). Per
our conventions: the wire fields exist in the canonical proto files and
changelog (v0.5.0, 2026-06-30) but in no tagged release --- the latest
tag is v0.4.1 (\texttt{research/upstream-defects.md}, item~3); the pool
and its activation are \emph{specified} (ZIP~2005, Proposed; ZIP~258,
Draft: Testnet activation $4{,}134{,}000$, Mainnet scheduled in deployed
node software at $3{,}428{,}143$); nothing deployed serves the
compact-format fields.

\subsection{Divergences between ZIP-307 and the deployed protocol}%
\label{sec:sync-compact-divergences}

Per our conventions, we flag rather than resolve the following.

\begin{itemize}
\item \emph{No normative specification of the deployed wire format.}
  ZIP~307 is Draft, Sapling-only, and its message definitions are stale ---
  old names, old field numbers, no Orchard, no \texttt{ChainMetadata}, no
  \texttt{vin}/\texttt{vout} --- while the \textsf{lightwallet-protocol}
  README disclaims being a specification and defers to ZIP~307
  (\texttt{zip-0307.rst}, ``Compact Stream Format'';
  \texttt{lightwallet-protocol/README.md}). The deployed format is
  specified only by tagged proto releases.
\item \emph{The grace-period pseudocode bug.} Both ``[Canopy onward]''
  conditions in ZIP~307's compact-decryption procedure test $\mathsf{height}
  < \mathsf{CanopyActivationHeight} + \mathsf{ZIP212GracePeriod}$; read
  literally they reject $\mathtt{0x01}$ during the grace period and
  constrain nothing after it. Deployed \texttt{zip212\_enforcement} accepts
  both lead bytes during the grace period and requires $\mathtt{0x02}$
  after (\texttt{zip-0307.rst}, lines~275--276;
  \textsf{zcash\_primitives},
  \texttt{src/transaction/}\allowbreak\texttt{components/}\allowbreak%
  \texttt{sapling.rs}, lines~27--40;
  \texttt{research/}\allowbreak\texttt{upstream-defects.md}, item~1,
  verified 2026-07-15 at \textsf{zips} \texttt{69610984}).
\item \emph{Inclusion semantics.} ZIP~307 says Sapling-relevant
  transactions only; the deployed default stream includes Orchard data,
  transparent data is available on request
  (\texttt{BlockRange.poolTypes}), and \texttt{GetBlock} returns all
  pools unconditionally (\S\ref{sec:sync-compact-servers}).
\item \emph{Three behaviours for \texttt{protoVersion}.}
  Server \textsf{lightwalletd} never sets it, \textsf{Zaino} sets~1, and
  v0.5.0 removed it with field number and name reserved
  (issue \texttt{zcash/lightwallet-protocol\#25}); \textsf{librustzcash}
  never reads it (\S\ref{sec:sync-compact-wire}).
\item \emph{Dead header-validation surface.} The proto declares
  \texttt{header} ``should always be unset''; ZIP~307's SPV-style header and
  \texttt{finalSaplingRoot} checks are unimplementable against deployed
  servers, and only \texttt{prevHash}/height continuity is enforced
  client-side (\S\ref{sec:sync-compact-wire},
  \S\ref{sec:sync-compact-scan}).
\item \emph{An advertised but unpopulated \texttt{fee} field.} The proto
  says ``present if server can provide''; \textsf{lightwalletd} has a TODO,
  \textsf{Zaino} hardcodes zero (\S\ref{sec:sync-compact-wire}).
\item \emph{Vendored-revision skew.} \textsf{lightwalletd} vendors v0.4.1,
  \textsf{librustzcash} main v0.5.0, and \textsf{Zaino} an intermediate
  snapshot with both the Ironwood fields and \texttt{protoVersion};
  \textsf{Zaino}'s vendored changelog dates v0.5.0 as 2026-07-03 without
  the \texttt{protoVersion} removal, while the canonical changelog dates it
  2026-06-30 with the removal --- and \textsf{Zaino} still sets a field the
  canonical proto reserves (diff of the three
  \texttt{compact\_formats.proto} copies;
  \texttt{zaino-proto/lightwallet-protocol/CHANGELOG.md}, lines~40--51).
\end{itemize}

%% Sync Guide, section: The light-client service. Body only; assumes the
%% shared preamble of the series (orchard-guide.tex lines 1-90).
%%
%% NOTE (draft-stage): the research memo counted "22 RPCs"; the vendored
%% walletrpc/service.proto at lightwalletd 61fee32e defines exactly 20
%% (verified firsthand: grep 'rpc ' service.proto, lines 223-325). The text
%% below uses the verified count of twenty.

\section{The light-client service}\label{sec:sync-service}

A light wallet holds keys and no chain. The \emph{Orchard Guide}
(\S``Transmission and detection of a note: encryption, trial decryption, and
synchronisation'') establishes what synchronisation consumes: every compact
shielded output on the chain, for trial decryption; every nullifier, for
spentness; the note commitment tree's frontier and subtree roots, for
witness assembly; and, for the few transactions that prove to be the
wallet's own, the full serialised bytes. The deployed layer that delivers
these is a single gRPC service, \texttt{cash.z.wallet.sdk.rpc.\allowbreak
CompactTxStreamer}, spoken between the wallet and an untrusted proxy in
front of a full node. This section documents the service as deployed: the
authority trail from ZIP~307 through the canonical protobuf definitions
(\S\ref{sec:sync-authority}), the compact data format
(\S\ref{sec:sync-compact}), the interaction model the ZIP prescribes
(\S\ref{sec:sync-interaction}), the twenty RPCs and their semantics
(\S\ref{sec:sync-rpcs}), the two server implementations --- lightwalletd,
the deployed proxy (\S\ref{sec:sync-lwd}), and Zaino, its successor
(\S\ref{sec:sync-zaino}) --- and a register of the points where the
implementations and the normative texts disagree
(\S\ref{sec:sync-divergences}). The linear costs this architecture accepts
--- trial decryption of every output, a maintained witness per note --- are
precisely the costs the Tachyon redesign removes (\emph{Tachyon Guide},
\S``What the fold removes'').

\subsection{Architecture and authority}\label{sec:sync-authority}

ZIP~307, ``Light Client Protocol for Payment Detection'' (Status: Draft),
fixes a three-component architecture: a Zcash full node as the root of
trust; a \emph{proxy server} that converts chain data into a compact
format; and the light client (ZIP~307, \S~``High-Level Design''). The
security model is deliberately narrow: the client obtains \emph{payment
detection privacy} against an honest-but-curious proxy, and trusts that
proxy to provide a correct view of the best chain; network-level privacy
(Tor, Nym) is assumed to be provided separately; spending privacy is out of
scope (ZIP~307, \S~``Security Model'').

\paragraph{Where the contract lives.} The wire contract is \emph{not} in
the ZIP. The two protobuf files that define the service ---
\texttt{service.proto} and \texttt{compact\_formats.proto} --- have their
canonical home in the \texttt{zcash/lightwallet-protocol} repository, which
self-describes as the canonical Protobuf definitions but explicitly not a
specification: ``The Zcash Light Client Protocol is described in ZIP-307''
(\textsf{lightwallet-protocol}, \texttt{README.md}). The authority trail is
therefore circular --- the ZIP defers detail to the deployed API, and the
API repository defers normative force back to the ZIP --- and the circle
does not close: ZIP~307's own protobuf sketch (\texttt{CompactSpend} and
\texttt{CompactOutput} messages, a \texttt{BlockID} embedded in
\texttt{CompactBlock}, a \texttt{RangeFilter}, a \texttt{FullTransaction}
return type, a four-RPC \texttt{CompactTxStreamer}) differs from the
deployed proto (\texttt{CompactSaplingSpend}, \texttt{CompactSaplingOutput},
\texttt{CompactOrchardAction}, flat height and hash fields,
\texttt{BlockRange}, \texttt{RawTransaction}, twenty RPCs), and the ZIP
warns that ``the details of the API described below may differ from the
implementation'' (ZIP~307, \S~``Compact Stream Format'', \S~``Proxy
operation''). We treat the canonical proto files, versioned and
change-logged, as the deployed contract, and cite the ZIP for architecture
and rationale. Both lightwalletd and Zaino vendor the proto files by git
subtree (\textsf{lightwallet-protocol}, \texttt{README.md}).

\paragraph{Contract version.} The canonical protocol is at v0.4.1
(2026-02-20). Version 0.4.0 (2025-12-03) added transparent data to the
compact format --- \texttt{PoolType}, \texttt{CompactTxIn}, \texttt{TxOut},
\texttt{CompactTx.vin}/\texttt{vout}, \texttt{BlockRange.poolTypes} --- and
the \texttt{LightdInfo} fields \texttt{upgradeName},
\texttt{upgradeHeight}, and \texttt{lightwalletProtocolVersion}, and
declared \texttt{GetBlockNullifiers} and
\texttt{GetBlockRangeNullifiers} deprecated; v0.4.1 designated the pair
deprecated in the proto itself (\texttt{option deprecated = true},
\texttt{service.proto:247,267}) and documented \texttt{GetBlock}'s
transparent-data behaviour (\textsf{lightwallet-protocol},
\texttt{CHANGELOG.md}).

\paragraph{The proto forks.} The vendored copies have already diverged
from the canonical file, in one direction: fields for the future Ironwood
pool. The canonical v0.4.1 has \texttt{PoolType} $=$ \{\texttt{INVALID},
\texttt{TRANSPARENT}, \texttt{SAPLING}, \texttt{ORCHARD}\} and
\texttt{ShieldedProtocol} $=$ \{\texttt{sapling}, \texttt{orchard}\}.
The fork in \textsf{librustzcash}
(\texttt{zcash\_client\_backend/proto/service.proto}) adds
\texttt{IRONWOOD = 4}, \texttt{TreeState.ironwoodTree = 7}, and
\texttt{ShieldedProtocol.ironwood = 2}. Zaino's fork
(\texttt{zaino-proto}) adds \texttt{IRONWOOD = 4},
\texttt{ironwoodTree = 7}, \texttt{ChainMetadata.\allowbreak
ironwoodCommitmentTreeSize = 3}, and \texttt{CompactTx.ironwoodActions =
9}, but not \texttt{ShieldedProtocol.ironwood}, and omits the
\texttt{deprecated} options on the two nullifier RPCs (diff of the three
proto trees; \textsf{zaino}, \texttt{docs/internal\_spec.md},
\S~``Zaino-Proto''). Zaino's documentation states the local copies exist
``for development purposes'' with a plan to rely on librustzcash's. The
``as deployed'' wire surface therefore differs by repository; this volume
cites the canonical file and flags fork-only fields where they appear.

\begin{remark}[Sources and pinning]\label{rem:sync-pins}
Deployed-behaviour claims in this section were verified against
\textsf{lightwalletd} master \texttt{61fee32e} (2026-06-25),
\textsf{zaino} \texttt{0e057e22} (2026-07-04), and \textsf{librustzcash}
\texttt{e30517e4} (2026-07-10); Zaino pins \textsf{zebra-chain} 11.0.0
from crates.io (\texttt{Cargo.lock}). File-and-line citations below refer
to these revisions.
\end{remark}

\subsection{The compact format}\label{sec:sync-service-compact}

The proxy's one transformation is compression: it strips from each
transaction everything a detecting wallet does not need. Proofs,
signatures, and value commitments verify the chain, and the client has
delegated chain verification to the proxy; what remains is exactly the
trial-decryption input and the nullifier.

\begin{definition}[Compact Sapling output, ZIP~307]
\label{def:sync-compact-output}
A compact output retains three fields of a Sapling Output description:
\[
\mathtt{CompactSaplingOutput} \;=\;
\bigl(\mathtt{cmu}[32],\; \mathtt{ephemeralKey}[32],\;
\mathtt{ciphertext} = \mathtt{encCiphertext}[0..52]\bigr),
\]
116~bytes per output against $580 + 32$ bytes for the full description
with its commitment --- an ${\sim}80\%$ reduction (ZIP~307, \S~``Output
Compression''; \texttt{compact\_formats.proto:114--122}). The 52-byte
ciphertext prefix covers the compact note plaintext --- the contents and
opening of the note commitment. Truncation discards the AEAD tag;
integrity of the decrypted prefix is instead checked by recomputing the
note commitment, as the \emph{Orchard Guide} constructs in
\S``Transmission and detection of a note: encryption, trial decryption,
and synchronisation''.
\end{definition}

\begin{definition}[Compact Orchard action, deployed proto]
\label{def:sync-compact-action}
The deployed format extends the same compression to Orchard, one message
per Action:
\[
\mathtt{CompactOrchardAction} \;=\;
\bigl(\mathtt{nullifier}[32],\; \mathtt{cmx}[32],\;
\mathtt{ephemeralKey}[32],\; \mathtt{ciphertext}[52]\bigr)
\]
(\texttt{compact\_formats.proto:124--130}). The Orchard message carries
its nullifier inline because an Action both spends and creates; on the
Sapling side, spend compression is a separate message that reduces the
384-byte Spend description to its 32-byte nullifier --- a ${\sim}90\%$
reduction (ZIP~307, \S~``Spend Compression'').
\end{definition}

\paragraph{Transparent data.} Since v0.4.0 the compact format can carry
transparent inputs and outputs (\texttt{CompactTxIn}, \texttt{TxOut},
\texttt{CompactTx.vin}/\texttt{vout}); v0.4.1 fixes the convention that
the coinbase transaction's single null-outpoint \texttt{vin} is omitted,
and that clients identify the coinbase by testing
$\mathtt{CompactTx.index} = 0$ (\texttt{service.proto:225--237};
\texttt{compact\_formats.proto:70--76}).

\begin{remark}[The fee field is dead on the wire]\label{rem:sync-fee}
The proto documents \texttt{CompactTx.fee} as ``present if server can
provide'', with the stateless-server formula for the case of no
transparent inputs,
\[
\mathtt{fee} \;=\; \mathtt{valueBalanceSapling} +
\mathtt{valueBalanceOrchard} + \textstyle\sum \mathtt{vPubNew} -
\sum \mathtt{vPubOld} - \sum \mathtt{tOut}
\]
(\texttt{compact\_formats.proto:59--64}); with transparent inputs a
stateless server cannot compute the fee at all. Neither deployed server
ever sets the field: lightwalletd's parser leaves it at~0 with a
``TODO: calculate fees'' comment (\textsf{lightwalletd},
\texttt{parser/transaction.go:397--432}), and Zaino's conversions pass
\texttt{fee = None}, encoded~0 (\textsf{zaino},
\texttt{zaino-state/src/chain\_index/}\allowbreak
\texttt{types/db/legacy.rs:1397--1470}).
The deployed contract is therefore that the field carries no information,
and a client must treat it as absent; we flag the proto's ``present if''
language as unimplemented on both servers.
\end{remark}

\subsection{The interaction model of ZIP~307}\label{sec:sync-interaction}

The ZIP prescribes a four-phase client-server interaction (ZIP~307,
\S~``Client-server interaction''):

\begin{enumerate}[label=(\Alph*)]
\item \emph{Startup.} The client fetches the latest \texttt{BlockID}. A
  wallet importing a pre-existing seed with no recorded birthday starts at
  Sapling activation, since no shielded address it can derive predates
  Sapling (ZIP~307, \S~``Importing a pre-existing seed'').
\item \emph{First sync.} The client requests
  $\mathtt{GetBlockRange}(X, Y)$ and, on any subsequent request, verifies
  the hash of the overlap block --- the last block it already holds ---
  against its cached copy, so that a reorganisation is detected as a hash
  mismatch.
\item \emph{Transaction fetching.} Detected notes are completed by
  fetching full transactions via \texttt{GetTransaction}. Each such fetch
  names a txid to the proxy, so the ZIP mandates cover traffic: clients
  SHOULD download decoy transactions, or all transactions of a touched
  block, and MUST warn users of the privacy loss otherwise (ZIP~307,
  \S~``Transaction privacy'').
\item \emph{Incremental sync.} Later syncs repeat (B) from the cached
  tip, re-checking the overlap block.
\end{enumerate}

Two hardening measures described by the ZIP are not deployed. Block-header
validation by the client is described as only partially implemented ---
only \texttt{prevHash} linkage is checked (ZIP~307, \S~``Block header
validation''). The \emph{bucketing} privacy enhancement --- for bucket
size $N$, start each range request at
\[
\mathrm{floor}(X) \;=\; X - (X \bmod N)
\]
instead of at the true resume height $X$, hash-checking and discarding
blocks in $[\mathrm{floor}(X), X]$, so that the request leaks only the
bucket and absorbs reorganisations --- is proposed and not implemented
(ZIP~307, \S~``Block privacy via bucketing''). Both are ZIP-versus-deployed
gaps a reader should not assume closed.

The reference client narrows the surface further: the sync loop in
\textsf{zcash\_client\_backend}
(\texttt{src/sync.rs:1--10,244--381}) exercises four shielded-sync RPCs
(below) plus, when built with feature \texttt{transparent-inputs},
\texttt{GetAddressUtxosStream}, issued per account per cycle by
\texttt{refresh\_utxos} (\texttt{src/sync.rs:117--126,511--527}; cf.\
\S\ref{sec:sync-privacy-schedule}):
\begin{itemize}
\item \texttt{get\_latest\_block}, in \texttt{update\_chain\_tip};
\item \texttt{get\_subtree\_roots}, in \texttt{download\_subtree\_roots};
\item \texttt{get\_block\_range}, in \texttt{download\_blocks};
\item \texttt{get\_tree\_state}, in \texttt{download\_chain\_state}.
\end{itemize}
Transaction enhancement via \texttt{GetTransaction} is listed as not yet
implemented in that module (\texttt{src/sync.rs}). We return to the loop
in the next section.
%% TODO: replace the prose forward reference with \ref{sec:sync-loop} (or
%% the final label of the scanning-pipeline section) once it is drafted.

\subsection{The RPC surface}\label{sec:sync-rpcs}

The service defines twenty RPCs (\texttt{service.proto:221--326});
Table~\ref{tab:sync-rpcs} groups them. Three are deprecated, one is
testing-only, and six serve the transparent layer, outside this volume's
shielded-sync scope. The four marked \checkmark\ are the ones the
reference sync loop consumes for shielded sync
(\S\ref{sec:sync-interaction}).

\begin{table}[htbp]
\centering
\small
\begin{tabular}{@{}llcl@{}}
\toprule
Group & RPC & Sync loop & Notes \\
\midrule
Chain & \texttt{GetLatestBlock} & \checkmark & tip height and hash \\
 & \texttt{GetBlock} & & one block, all pools \\
 & \texttt{GetBlockRange} & \checkmark & streamed range, pool filter \\
 & \texttt{GetBlockNullifiers} & & deprecated (v0.4.0; proto option
   v0.4.1) \\
 & \texttt{GetBlockRangeNullifiers} & & deprecated (v0.4.0; proto option
   v0.4.1) \\
\addlinespace
Trees & \texttt{GetTreeState} & \checkmark & frontiers at height/hash \\
 & \texttt{GetLatestTreeState} & & frontiers at the tip \\
 & \texttt{GetSubtreeRoots} & \checkmark & depth-16 subtree roots \\
\addlinespace
Transactions & \texttt{GetTransaction} & & full bytes by txid \\
 & \texttt{SendTransaction} & & submit raw bytes \\
\addlinespace
Mempool & \texttt{GetMempoolTx} & & compact mempool contents \\
 & \texttt{GetMempoolStream} & & full txs until next block \\
\addlinespace
Transparent & \texttt{GetTaddressTxids} & & deprecated, misnamed \\
 & \texttt{GetTaddressTransactions} & & \\
 & \texttt{GetTaddressBalance} & & \\
 & \texttt{GetTaddressBalanceStream} & & \\
 & \texttt{GetAddressUtxos} & & \\
 & \texttt{GetAddressUtxosStream} & & \\
\addlinespace
Metadata & \texttt{GetLightdInfo} & & server and chain metadata \\
 & \texttt{Ping} & & testing-only \\
\bottomrule
\end{tabular}
\caption{The twenty RPCs of \texttt{CompactTxStreamer}
(\texttt{service.proto:221--326}, lightwallet-protocol v0.4.1). The
\checkmark\ column marks the shielded-sync RPCs exercised by the
reference sync loop; with feature \texttt{transparent-inputs} the loop
also issues \texttt{GetAddressUtxosStream} per account per cycle
(\textsf{zcash\_client\_backend}, \texttt{src/sync.rs}).}
\label{tab:sync-rpcs}
\end{table}

\begin{remark}[Hash orders]\label{rem:sync-endian}
Two byte orders cross the interface. The byte-typed fields
\texttt{BlockID.hash} and \texttt{CompactBlock.hash} carry raw bytes in
little-endian (protocol) order; lightwalletd reverses the
big-endian hex it receives from the node's RPCs
(\textsf{lightwalletd}, \texttt{frontend/service.go:69--90}). Field
\texttt{TreeState.hash}, by contrast, is big-endian display hex
\emph{text} (\texttt{service.proto:176--184,307--312}). Requests follow
the same split: \texttt{TxFilter} takes protocol-order txid bytes that
the server reverses into display hex for \texttt{getrawtransaction}
(\texttt{frontend/service.go:403--442}), and the mempool exclude entries
are protocol-order byte suffixes the server reverses into display-order
prefixes (\texttt{frontend/service.go:600--747}). Every byte-typed hash
on this interface is protocol order; every string-typed hash is display
order.
\end{remark}

\paragraph{Chain tip.} The RPC
$\mathtt{GetLatestBlock}(\mathtt{ChainSpec})$ returns a \texttt{BlockID}
--- height and hash of the best chain tip; lightwalletd proxies
\texttt{getblockchaininfo} and returns the hash as little-endian raw bytes
(\textsf{lightwalletd}, \texttt{frontend/service.go:69--90}). This is the
chain tip the wallet's confirmation accounting consumes (\emph{Wallet
Guide}, \S``Confirmations, trust, and spendability'').

\paragraph{Single block.} The RPC $\mathtt{GetBlock}(\mathtt{BlockID})$
returns one \texttt{CompactBlock}, and per the v0.4.1 contract includes
transaction data for \emph{all} value pools, transparent
\texttt{vin}/\texttt{vout} included --- unlike \texttt{GetBlockRange},
which defaults to shielded-only (\texttt{service.proto:225--237}).
The proxy lightwalletd serves the request from its block cache first and
falls back to the node on a miss; a height above the tip returns gRPC
\texttt{OutOfRange} (``block \%d is newer than the latest block''), a
lookup by hash returns \texttt{InvalidArgument} (``not yet implemented'',
referencing PR~\#309), and height~0 with no hash is
\texttt{InvalidArgument} (\textsf{lightwalletd},
\texttt{common/common.go:495--624}, \texttt{frontend/service.go:166--189}).
Zaino accepts lookup by hash --- hash takes precedence and is resolved to
a height via its chain index --- and builds the block with all pools
included, matching the v0.4.1 contract; the proto makes hash support
optional, so both behaviours are conformant, but they diverge
(\textsf{zaino}, \texttt{zaino-state/src/backends/state.rs:1954--2058};
\texttt{zaino-proto/src/proto/utils.rs:310--323}).

\paragraph{Block ranges.} The RPC
$\mathtt{GetBlockRange}(\mathtt{BlockRange})$ streams consecutive
\texttt{CompactBlock}s over the closed interval
$[\min(s,e), \max(s,e)]$, in increasing height order iff $s \le e$ and
decreasing otherwise --- lightwalletd computes
$j = \mathit{high} - (i - \mathit{low})$, Zaino tests
$s \le e$ (\texttt{service.proto:250--255};
\textsf{lightwalletd}, \texttt{common/common.go:574--624};
\textsf{zaino}, \texttt{zaino-state/}\allowbreak\texttt{src/}\allowbreak
\texttt{chain\_index.rs:1774--1802}). An
empty \texttt{poolTypes} field selects the legacy behaviour: prune each
block to shielded (Sapling and Orchard) data only; clients MUST verify
server capability before setting \texttt{poolTypes} non-empty
(\texttt{service.proto:28--42}). Filtering prunes transaction components
and then drops empty transactions, but a block whose transactions were all
filtered is still returned, with empty \texttt{vtx}
(\textsf{lightwalletd}, \texttt{common/common.go:574--624}). Zaino
validates the request up front --- start and end must be present and fit
in \texttt{u32}; \texttt{POOL\_TYPE\_INVALID} or unknown pool values are
\texttt{InvalidArgument} --- streams through an mpsc channel of the
configured \texttt{channel\_size}, caps an ascending range at the tip and
appends a trailing \texttt{OutOfRange} error after the valid blocks, and
bounds production by a hard timeout of $4 \times$ the service timeout
(\textsf{zaino}, \texttt{zaino-state/src/backends/state.rs:550--683};
\texttt{zaino-proto/src/proto/utils.rs:53--152}).

\paragraph{Nullifier variants.} The RPCs \texttt{GetBlockNullifiers}
and \texttt{GetBlockRangeNullifiers} ---
both deprecated (v0.4.0 changelog; proto option v0.4.1) in favour of
\texttt{GetBlockRange} with
\texttt{poolTypes} --- return \texttt{CompactBlock}s reduced to shielded
nullifier data: lightwalletd strips Sapling outputs and transparent
\texttt{vin}/\texttt{vout}, reduces each Orchard action to its nullifier,
and zeroes both commitment-tree sizes; \texttt{GetBlockRangeNullifiers}
MUST ignore any \texttt{TRANSPARENT} member of \texttt{poolTypes}, and
lightwalletd filters it out (\texttt{service.proto:239--268};
\textsf{lightwalletd}, \texttt{frontend/service.go:191--225,259--305}).
Zaino additionally accepts these by hash, as for \texttt{GetBlock}
(\textsf{zaino}, \texttt{zaino-state/src/backends/fetch.rs:929--1041}).

\paragraph{Tree state.} The RPC
$\mathtt{GetTreeState}(\mathtt{BlockID})$ returns the note commitment tree
state for a block specified by height or hash: network name, height, block
hash as a hex string, block time, and the \texttt{saplingTree} and
\texttt{orchardTree} frontiers as hex (\texttt{service.proto:176--184}).
In lightwalletd the reply derives from the node's
\texttt{z\_gettreestate} RPC,
looping along \texttt{Sapling.SkipHash} when a reply carries no
\texttt{FinalState}; \texttt{GetLatestTreeState} calls the same path at
\texttt{getblockchaininfo}'s tip height (\textsf{lightwalletd},
\texttt{frontend/service.go:307--401}). The frontier is what seeds the
wallet's shard tree below its birthday, so that witness assembly never
replays the chain's prefix (\emph{Orchard Guide}, \S``Proof of ownership
of \emph{some} valid note: the commitment tree''). Zaino resolves
hash-or-height, reports the BIP70 network name, and returns a seventh
field \texttt{ironwoodTree} (hex, empty when absent) that exists only in
Zaino's and librustzcash's proto forks (\S\ref{sec:sync-authority});
\texttt{GetLatestTreeState} reuses \texttt{GetTreeState} at the current
chain height (\textsf{zaino},
\texttt{zaino-state/src/backends/state.rs:2516--2576};
\texttt{zaino-proto/proto/service.proto:183}). One stale cross-reference:
the proto comment for \texttt{GetTreeState} cites ``section 3.7 of the
Zcash protocol specification'', while the current
\texttt{protocol.tex} numbers note commitment trees as \S~3.8.

\paragraph{Subtree roots.} Shielded outputs of a pool are numbered from
genesis, and each consecutive group of $2^{16} = 65536$ note commitments
forms one depth-16 subtree of the pool's commitment tree; the stream
element is
\begin{align*}
\mathtt{SubtreeRoot}_i \;=\; \bigl(&\text{Merkle root of subtree } i,\\
&\text{hash and height of the block containing output }
(i+1)\cdot 2^{16} - 1\bigr),
\end{align*}
so that, for example, Sapling output $65535$ --- the last of subtree $0$
--- falls in block $558822$ (\texttt{service.proto:191--200};
\textsf{lightwalletd}, \texttt{common/common.go:178--219}). The RPC
$\mathtt{GetSubtreeRoots}(\mathtt{GetSubtreeRootsArg})$ streams these for
the chosen shielded protocol from \texttt{startIndex}, with
$\mathtt{maxEntries} = 0$ meaning all. The lightwalletd handler calls
\texttt{z\_getsubtreesbyindex} and then, per subtree, fetches the
completing block via \texttt{GetBlock} at its end height to obtain the
hash, reversed to big-endian display-order bytes, the opposite
convention from \texttt{CompactBlock.hash}
(Remark~\ref{rem:sync-tree-byteorder}; \textsf{lightwalletd},
\texttt{frontend/service.go:832--907}). These roots are exactly the
retained internal nodes from which the wallet's shard tree assembles
authentication paths without replaying every commitment (\emph{Orchard
Guide}, \S``Proof of ownership of \emph{some} valid note: the commitment
tree''). Zaino converts \texttt{startIndex} and \texttt{maxEntries} to
\texttt{u16} and returns \texttt{InvalidArgument} when either exceeds
\texttt{u16} range, where lightwalletd forwards the full \texttt{uint32}
--- bounded by design, since a depth-32 tree has at most $2^{16}$
depth-16 subtrees, but the wire contracts differ; it fills the completing
block data via \texttt{z\_get\_block} per subtree, under a $4 \times$
service-timeout deadline. One Zebra-specific caveat: during Zebra's
initial subtree-index rebuild the underlying RPC returns an empty list
for not-yet-rebuilt indices, where zcashd rebuilds before serving
(\textsf{zaino}, \texttt{zaino-state/src/indexer.rs:429--452,754--912}).

\paragraph{Full transactions.} The RPC
$\mathtt{GetTransaction}(\mathtt{TxFilter})$ returns the full serialised
transaction. Only lookup by 32-byte txid is supported; a block-plus-index
lookup returns \texttt{InvalidArgument} (``specify a txid, not a
blockhash+num''). The lightwalletd handler converts the little-endian
txid bytes to big-endian hex, calls \texttt{getrawtransaction} with
\texttt{verbose=1}, and maps RPC failure to gRPC \texttt{NotFound}
(\texttt{service.proto:44--51,270--271}; \textsf{lightwalletd},
\texttt{frontend/service.go:403--442}). The privacy cost of this call ---
it names a txid of interest to the proxy --- is the reason for the ZIP's
decoy mandate (\S\ref{sec:sync-interaction}).

\paragraph{Submission.} The RPC
$\mathtt{SendTransaction}(\mathtt{RawTransaction})$ submits via the
node's \texttt{sendrawtransaction}. Success returns
$\mathtt{SendResponse}\{\mathtt{errorCode}{=}0,\
\mathtt{errorMessage}{=}\text{txid hex}\}$; a node rejection is parsed
from the ``code: message'' error string into a non-zero
\texttt{errorCode} plus message --- node-level rejection travels in-band,
not as a gRPC error, so a transport-level success is not an acceptance
(\texttt{service.proto:83--89}; \textsf{lightwalletd},
\texttt{frontend/service.go:454--509}). Wallet-side obligations around
this call --- store-then-broadcast ordering, rebroadcast, expiry --- are
the \emph{Wallet Guide}'s subject (\S``Store, then broadcast'').

\begin{remark}[The \texttt{RawTransaction.height} sentinels]
\label{rem:sync-height-sentinel}
An original protobuf-definition error typed the field \texttt{uint64}
rather than \texttt{int64}, forcing sentinel semantics
(\texttt{service.proto:53--81}; \textsf{lightwalletd},
\texttt{common/common.go:626--661}):
\[
\mathtt{height} \;=\;
\begin{cases}
0 & \text{transaction in the mempool,}\\
2^{64}-1 \; (\mathtt{0xffffffffffffffff}) &
  \text{mined on a non-main-chain fork (zcashd's } -1\text{),}\\
h & \text{mined on the main chain at height } h.
\end{cases}
\]
The proto's own comments contradict each other: the message-level comment
says the field carries ``the latest block height'' when returned by
\texttt{GetMempoolStream}, the field-level comment specifies the sentinel
mapping above, and a FIXME cites librustzcash issue~\#1484. The deployed
servers realise both readings --- lightwalletd streams $0$, Zaino streams
the current best-tip height (\S\ref{sec:sync-divergences}). We flag the
mempool-height contract as unresolved
(\texttt{research/upstream-defects.md}, which tracks this as the fourth
known defect, already filed as zcash/librustzcash\#1484).
\end{remark}

\paragraph{Mempool contents.} The RPC
$\mathtt{GetMempoolTx}(\mathtt{GetMempoolTxRequest})$ streams
\texttt{CompactTx} for the current mempool contents, possibly a few
seconds stale. The request may carry \texttt{exclude\_txid\_suffixes}:
txid byte-suffixes of any length up to 32 bytes, in client-side protocol
order; the server reverses and hex-encodes each into a display-order
prefix, and a mempool transaction is excluded iff its txid matches a
prefix that matches \emph{exactly one} mempool txid --- when two or more
txids match, none of the matching transactions are excluded, and
unmatched entries are ignored (\texttt{service.proto:157--174,289--301};
\textsf{lightwalletd}, \texttt{frontend/service.go:600--747}). Field~2 of
the request is reserved for a future \texttt{exclude\_txid\_prefixes}.
An empty \texttt{poolTypes} defaults to Sapling-plus-Orchard pruning;
\texttt{POOL\_TYPE\_INVALID} is rejected. The lightwalletd
implementation backs the RPC with
a process-global cache --- \texttt{mempoolMap} from txid to
\texttt{CompactTx}, plus \texttt{mempoolList} --- refreshed at most every
2~s from \texttt{getrawmempool} plus per-transaction
\texttt{getrawtransaction} \texttt{verbose=0}, reusing already-fetched
entries across refreshes; mempool \texttt{CompactTx}s are built with
height~0 and no fee (\textsf{lightwalletd},
\texttt{frontend/service.go:589--708}). Zaino enforces the same 32-byte
suffix bound and the same exactly-one-match exclusion rule, parses
transactions from stored serialised bytes, and prunes by the requested
pool types before streaming (\textsf{zaino},
\texttt{zaino-state/src/backends/state.rs:2302--2429},
\texttt{chain\_index/mempool.rs:381--417}).

\paragraph{Mempool stream.} The RPC
$\mathtt{GetMempoolStream}(\mathtt{Empty})$ streams full
\texttt{RawTransaction}s of the current mempool and keeps the stream
open, closing it when a new block is mined
(\texttt{service.proto:303--305}). The lightwalletd implementation is
process-global state --- \texttt{g\_txList} in arrival order, a
\texttt{g\_txidSeen} set, \texttt{g\_lastBlockChainInfo} --- polled at
most every 2~s via \texttt{getblockchaininfo} and
\texttt{getrawmempool}; the first stream to observe a tip-hash change
resets the shared state and all open streams terminate; each per-client
loop sleeps 200~ms between sends; transactions mined since listing
(height $\neq 0$) are skipped; streamed transactions carry
$\mathtt{height} = 0$ (\textsf{lightwalletd},
\texttt{common/mempool.go:14--152},
\texttt{frontend/service.go:581--587}). Zaino streams
$\mathtt{height} = {}$current best-tip height, closes the stream when its
mempool status flips to \texttt{Syncing} on a tip change, and
additionally imposes a hard $6 \times$ service-timeout deadline ---
180~s at defaults --- where lightwalletd keeps the stream open until the
next block regardless of duration; a long block interval can therefore
end a Zaino stream with a deadline error where the proto's documented
close-on-new-block semantics would hold (\textsf{zaino},
\texttt{zaino-state/src/backends/state.rs:2433--2510},
\texttt{chain\_index/mempool.rs:419--480}).

\paragraph{Server metadata.} The RPC
$\mathtt{GetLightdInfo}(\mathtt{Empty})$ returns server metadata plus
chain state assembled from \texttt{getinfo} and
\texttt{getblockchaininfo}: version, vendor (``ECC LightWalletD''),
\texttt{taddrSupport} true, chain name, the Sapling activation height ---
looked up via the Sapling consensus branch ID $\mathtt{0x76b809bb}$ in
the upgrades map --- the chain tip's consensus branch ID, block height,
build and git fields, \texttt{estimatedHeight} (less than the tip while
the node itself syncs), donation address, and the next pending upgrade's
name and height (\texttt{service.proto:97--118}; \textsf{lightwalletd},
\texttt{common/common.go:261--319}). The proxy does \emph{not}
populate the proto's \texttt{lightwalletProtocolVersion} field~18. Zaino
reports vendor ``ZingoLabs ZainoD'', hard-codes
$\mathtt{lightwallet\_protocol\_version} = \text{``v0.5.0''}$ --- a
version with no tagged release in the canonical repository, the latest
tag being v0.4.1, though an untagged v0.5.0 changelog section
(2026-06-30) exists (\texttt{research/upstream-defects.md}, item~3;
\texttt{zaino-state/src/backends/fetch.rs:2097},
\texttt{state.rs:2776}) --- defaults the Sapling activation height to 1
on regtest, and
carries the connected Zebra node's build information in the
\texttt{zcashdBuild}/\texttt{zcashdSubversion} fields (\textsf{zaino},
\texttt{zaino-state/src/backends/state.rs:2724--2793}). Both sides are
flagged in \S\ref{sec:sync-divergences}.

\paragraph{Ping.} The RPC \texttt{Ping} is testing-only; lightwalletd
disables it unless started with the flag
\texttt{-{}-ping-very-insecure}, because the call lets a client create
arbitrarily many concurrent server threads --- a
resource-exhaustion risk (\texttt{service.proto:324--325};
\textsf{lightwalletd}, \texttt{frontend/service.go:920--936}). Zaino
returns gRPC \texttt{Unimplemented} (\textsf{zaino},
\texttt{zaino-state/src/backends/state.rs:2782--2792}).

\paragraph{Transparent-address RPCs.} The six remaining RPCs
(\texttt{GetTaddressTxids} --- deprecated and misnamed, it returns
transactions --- \texttt{GetTaddressTransactions},
\texttt{GetTaddressBalance}, \texttt{GetTaddressBalanceStream},
\texttt{GetAddressUtxos}, \texttt{GetAddressUtxosStream}) proxy the
node's insight-explorer-derived RPCs \texttt{getaddresstxids},
\texttt{getaddressbalance}, and \texttt{getaddressutxos}; they serve the
transparent layer and fall outside this volume's shielded-sync scope.
For these, lightwalletd validates t-addresses against the regex
\texttt{\textbackslash At[a-zA-Z0-9]\{34\}\textbackslash z} and caps each
\texttt{GetTaddressTransactions} inner fetch loop with a 30-second
context timeout (\textsf{lightwalletd},
\texttt{frontend/service.go:55--164,511--579,749--830}). Zaino caps the
request's address list at $\mathtt{UTXO\_MAX\_ADDRESSES} = 1000$
(\texttt{InvalidArgument} beyond) --- a server-side fan-out bound
lightwalletd lacks --- and honours \texttt{start\_height} and
\texttt{max\_entries} ($0$ meaning unlimited) as response-size bounds
(\textsf{zaino}, \texttt{zaino-state/src/backends.rs:29--45},
\texttt{backends/state.rs:2591--2719}).

\subsection{lightwalletd: the deployed proxy}\label{sec:sync-lwd}

The deployed server is a Go proxy: each gRPC handler in
\texttt{frontend/service.go} translates its request into node JSON-RPC
calls through the indirect function \texttt{common.RawRequest}. It
supports both zcashd and zebrad backends, detecting the backend from the
node's subversion string (\texttt{/Zebra:} versus \texttt{/MagicBean:});
a zcashd backend must have the \texttt{lightwalletd} or
\texttt{insightexplorer} experimental feature enabled or lightwalletd
exits; the default assumed node is zebrad (\textsf{lightwalletd},
\texttt{cmd/root.go:200--263}, \texttt{common/common.go:27--35,61--64}).

\paragraph{Ingestion.} A single \texttt{BlockIngestor} goroutine polls
\texttt{getbestblockhash} every 2~s; when the tip differs it fetches the
next block via \emph{two} \texttt{getblock} calls --- first
\texttt{verbose=1} for the txids and hash, a workaround because the
built-in parser miscomputes v5 transaction IDs (issue~\#392), then the
raw block by that hash, so a reorganisation between the calls cannot mix
two blocks. A \texttt{prevHash} mismatch against the cache tip is treated
as a reorganisation and exactly one block is dropped
(\texttt{Reorg(height-1)}) per iteration; a \texttt{getblock} failure
retries after 8~s (\textsf{lightwalletd},
\texttt{common/common.go:321--493}). The 120~s sleep
(\texttt{common/common.go:483--489}) runs only when the cache is empty
($\mathrm{height} = \texttt{GetFirstHeight()}$) and the node returned no
block at that height; with the cache now starting at height 0
(\texttt{cmd/root.go:296--299}) the branch is reachable only against a
node that cannot yet serve genesis, and its log line (``Waiting for
\ldots\ Sapling activation height'') is stale wording, not stale
behaviour.

\paragraph{The block cache.} The only persistent state lightwalletd
maintains is the compact-block cache: two append-only files per chain,
\texttt{<data-dir>/db/<chain>/lengths} and \texttt{.../blocks}. With
$d = \mathrm{protobuf}(\mathtt{CompactBlock})$ the entry for height $h$
is
\[
\mathrm{FNV1a}_{64}\bigl(\mathrm{LE}_{64}(h) \,\|\, d\bigr) \,\|\, d
\quad\text{in \texttt{blocks}},\qquad
\mathrm{LE}_{32}\bigl(|d|\bigr) \quad\text{in \texttt{lengths}},
\]
with the offset table reconstructed on startup and valid lengths
constrained to $74 \le |d| \le 4{,}000{,}000$ bytes; corruption --- a bad
checksum, an impossible length, a partial entry --- triggers automatic
cache clearing and redownload (\textsf{lightwalletd},
\texttt{common/cache.go:21--230}). The cache now starts at height~0,
previously Sapling activation, because compact blocks carry transparent
data (\textsf{lightwalletd}, \texttt{cmd/root.go:265--300}). All other
responses are proxied per-request; the mempool caches of
\S\ref{sec:sync-rpcs} are in-memory only; no per-client state is kept.

\paragraph{Operational defaults.} gRPC binds 127.0.0.1:9067 and HTTP
127.0.0.1:9068 (Prometheus \texttt{/metrics}); TLS is required --- the
server exits if the certificate or key is unloadable --- unless
\texttt{-{}-no-tls-very-insecure} (plaintext) or
\texttt{-{}-gen-cert-very-insecure} (self-signed); the data directory
defaults to \texttt{/var/lib/lightwalletd}; \texttt{-{}-nocache} disables
the disk cache; \texttt{-{}-sync-from-height} and
\texttt{-{}-redownload} control cache rebuild; the darkside mock mode
(\texttt{-{}-darkside-very-insecure}) auto-terminates after 30 minutes
(\textsf{lightwalletd}, \texttt{cmd/root.go:139--180,375--424,498--499}).

\subsection{Zaino: the successor indexer}\label{sec:sync-zaino}

Zaino is the successor indexer, in Rust: \texttt{zainod} is the daemon,
\texttt{zaino-serve} implements \texttt{CompactTxStreamer} over tonic,
\texttt{zaino-state} holds the chain index and backends,
\texttt{zaino-fetch} is the JSON-RPC client, and \texttt{zaino-proto}
vendors the protos. The stated motivations are the zcashd deprecation, a
Rust stack alongside librustzcash and Zebra, and a clean
indexer/validator trust boundary: indexing functionality moves out of
the validator into Zaino, with Zebra exposing its
\texttt{ReadStateService} for direct read access (\textsf{zaino},
\texttt{README.md}, \S~``Motivations'', \S~``Project Structure'';
\texttt{docs/internal\_spec.md}).

\paragraph{Backends.} Configuration \texttt{backend = 'fetch' | 'state'}
selects one of two chain-fetch backends. Backend \texttt{StateService} is
backed
by Zebra's \texttt{ReadStateService} plus a \texttt{TrustedChainSync}
(\texttt{init\_read\_state\_with\_syncer}): it reads Zebra's state store
directly, and at startup blocks until the syncer's tip height \emph{and}
hash match the validator's JSON-RPC tip --- hash compared because height
alone is ambiguous during a reorganisation. Backend \texttt{FetchService}
is backed by the node's JSON-RPC interface (zcashd back-compatibility, or a
remote Zebra) and is marked ``\texttt{\#[deprecated]}: Will be eventually
replaced by \texttt{BlockchainSource}''. Both feed the same
\texttt{NodeBackedChainIndex} (\textsf{zaino},
\texttt{zaino-state/src/backends/state.rs:1,117--260},
\texttt{backends/fetch.rs:1,83--100}). The shipped example configuration
nonetheless selects \texttt{backend = 'fetch'} --- the deprecated backend
--- and the repository documentation does not say which backend is the
recommended production deployment
(\textsf{zaino}, \texttt{docs/example\_configs/zainod.toml}); we flag
the question as open.

\paragraph{The chain index.} Server-side state is the
\texttt{ChainIndex}, with three components (\textsf{zaino},
\texttt{docs/internal\_spec.md}, \S~``Zaino-State'';
\texttt{zaino-state/}\allowbreak\texttt{src/}\allowbreak
\texttt{chain\_index.rs:1--100}). A \texttt{Mempool}
--- a dashmap broadcast map keyed by txid, whose sync task polls
\texttt{get\_best\_block\_hash} and on a tip change clears the map and
notifies subscribers (\texttt{chain\_index/mempool.rs}). A
\texttt{NonFinalisedState} --- an in-memory \texttt{ArcSwap} snapshot of
the top blocks of all chains. A \texttt{FinalisedState} --- all blocks
below the seam --- served either by a versioned LMDB-backed persistent
database (v1) or, with \texttt{ephemeral\_finalised\_state}, by a
stateless passthrough to the validator; large syncs and version
migrations run in the background while serving continues
(\texttt{chain\_index/finalised\_state.rs:1--80}). The seam is
\[
\mathrm{floor}(t) \;=\; t - D_{\mathrm{op}}
\quad\text{(saturating at genesis)},
\]
where $t$ is the tip height and
$D_{\mathrm{op}} = \mathtt{OPERATIONAL\_NFS\_DEPTH} =
\mathtt{MAX\_BLOCK\_REORG\_HEIGHT} + 1 = 1001$ --- Zebra's reorganisation
limit is local node policy, not consensus; the finalised database serves
heights $\le \mathrm{floor}(t)$, the in-memory state serves heights above
it, and in-memory retention is hard-capped at
$\mathtt{MAX\_NFS\_DEPTH} = D_{\mathrm{op}} + 10 = 1011$ blocks below the
tip (\textsf{zaino}, \texttt{zaino-common/src/consensus.rs:16--17},
\texttt{chain\_index/non\_finalised\_state.rs:23--58}; \textsf{zebra},
\texttt{zebra-chain/src/parameters/constants.rs:30}).

\paragraph{Serving.} Each gRPC handler delegates to a
\texttt{LightWalletIndexer} trait implemented by both backend
subscribers; streaming responses are tokio mpsc channels wrapped in typed
streams (\textsf{zaino},
\texttt{zaino-serve/src/rpc/grpc/service.rs:159--320}). The gRPC server
may bind a public address only with TLS configured --- rejected at
startup otherwise --- and the default listen address is 127.0.0.1:8137;
Zaino additionally serves a zcashd-style JSON-RPC server, unencrypted and
restricted to loopback or private bind addresses by default
(\textsf{zaino}, \texttt{README.md}, \S~``Server network exposure'';
\texttt{zainod/src/config.rs:239}). Service defaults: timeout 30~s,
\texttt{channel\_size} 32; stream deadlines are $4\times$ the timeout
(120~s), and $6\times$ for \texttt{GetMempoolStream} (180~s)
(\textsf{zaino}, \texttt{zaino-common/src/config/service.rs:8--20}).
Until the non-finalised snapshot exists, \texttt{GetLatestBlock} and
\texttt{GetMempoolStream} return an \texttt{UnavailableNotSyncedEnough}
error where lightwalletd would proxy the node's answer immediately; a
code TODO says this ``probably shouldn't be an error'', so the startup
behaviour may change (\textsf{zaino},
\texttt{zaino-state/src/backends/state.rs:1940--1951}).

\paragraph{Compatibility.} Zaino's integration test suite includes a
\texttt{client\_rpcs} partition that compares its LightWallet gRPC
outputs against lightwalletd's, and the project has committed to serving
all Zcash JSON-RPC services required by non-validator clients,
superseding zcashd in that role (\textsf{zaino},
\texttt{docs/internal\_spec.md}, \S~``Zaino-Testutils and
Integration-Tests''; \texttt{docs/rpc\_api.md}).

\subsection{Divergence register}\label{sec:sync-divergences}

Table~\ref{tab:sync-divergences} collects the contract-level points at
which the two deployed servers differ from each other or from the
normative texts; each row is developed, with citations, in the paragraph
of \S\ref{sec:sync-rpcs}--\S\ref{sec:sync-zaino} that covers the RPC.
The ZIP-versus-deployment gaps --- the message sketch, partial header
validation, unimplemented bucketing --- are flagged in
\S\ref{sec:sync-authority} and \S\ref{sec:sync-interaction}; the proto
forks in \S\ref{sec:sync-authority}.

\begin{table}[htbp]
\centering
\small
\begin{tabular}{@{}p{0.34\linewidth}p{0.27\linewidth}p{0.31\linewidth}@{}}
\toprule
Contract point & lightwalletd & Zaino \\
\midrule
\texttt{GetBlock}(\texttt{Nullifiers}) by hash (optional in proto) &
\texttt{InvalidArgument}, ``not yet implemented'' (PR~\#309) &
supported; hash takes precedence \\
\addlinespace
Stream deadlines (proto: \texttt{GetMempoolStream} closes on new block) &
none; open until next block &
hard $4\times$/$6\times$ timeout (120~s/180~s at defaults) \\
\addlinespace
\texttt{RawTransaction.height} in \texttt{GetMempoolStream} (proto
comments contradict; issue \#1484) &
$0$ &
current best-tip height \\
\addlinespace
\texttt{GetSubtreeRoots} argument range (proto: \texttt{uint32}) &
forwards full range to node &
\texttt{InvalidArgument} above \texttt{u16::MAX} \\
\addlinespace
\texttt{CompactTx.fee} (proto: ``present if server can provide'') &
always 0 (parser TODO) &
always 0 (\texttt{fee = None}) \\
\addlinespace
\texttt{lightwalletProtocolVersion} (\texttt{LightdInfo} field 18,
v0.4.0) &
never populated &
hard-coded ``v0.5.0'', no tagged canonical release (latest tag
v0.4.1) \\
\addlinespace
Behaviour before server is synced &
proxies the node immediately &
\texttt{UnavailableNotSyncedEnough} on \texttt{GetLatestBlock},
\texttt{GetMempoolStream} \\
\addlinespace
\texttt{Ping} &
disabled unless \texttt{-{}-ping-very-insecure} &
\texttt{Unimplemented} \\
\addlinespace
Transparent-address fan-out &
unbounded address lists &
capped at 1000 addresses \\
\bottomrule
\end{tabular}
\caption{Deployed divergences across the \texttt{CompactTxStreamer}
contract. Citations for each row appear in the covering paragraphs of
\S\ref{sec:sync-rpcs}--\S\ref{sec:sync-zaino}.}
\label{tab:sync-divergences}
\end{table}

Two rows deserve emphasis. The mempool-height row is a live
specification bug: the proto's two comments cannot both be satisfied,
the deployed servers realise one reading each, and the FIXME in the
canonical file (librustzcash issue~\#1484) records the conflict without
resolving it; a client must therefore treat the height of a streamed
mempool transaction as carrying no portable information beyond ``not yet
mined''. The fee row is a silent one: nothing on the wire distinguishes
``fee is zero'' from ``fee not provided'', so the proto's conditional
language is, in deployed terms, a constant.

\FloatBarrier

\section{The scanning pipeline}\label{sec:sync-scanning}
%% Ground truth: librustzcash @ e30517e4 (2026-07-10); lightwalletd @
%% 61fee32e (2026-06-25); zaino @ 0e057e22 (2026-07-04); zips @ 69610984
%% (2026-07-09). Line references are against these commits.

A light wallet holds no chain. Everything it knows --- which notes it owns,
which of them are spent, and where each sits in the note commitment tree ---
it reconstructs by scanning a stream of \emph{compact blocks} supplied by an
indexing server. This section develops the deployed reconstruction pipeline
of \textsf{librustzcash}'s \textsf{zcash\_client\_backend} and
\textsf{zcash\_client\_sqlite} crates: the compact input format
(\S\ref{sec:sync-compact}), the account birthday that bounds the scan
(\S\ref{sec:sync-birthday}), the priority queue that orders it
(\S\ref{sec:sync-queue}), chain-tip tracking (\S\ref{sec:sync-tip}), the
per-range scanning machinery (\S\ref{sec:sync-scan}), persistence
(\S\ref{sec:sync-persist}), reorganisation recovery
(\S\ref{sec:sync-reorg}), and the reference loop that assembles the parts
(\S\ref{sec:sync-loop}).

The cost shape of the pipeline deserves stating first. For every block in a
suggested range, every compact output and action of every transaction is
trial-decrypted against every prepared incoming viewing key of every tracked
account --- two ZIP-32 scopes per pool per account
(\S\ref{sec:sync-scan}) --- so synchronisation work scales with global chain
traffic multiplied by tracked keys, and is independent of the wallet's own
activity; restore-from-seed replays the whole computation from the birthday
(\textsf{zcash\_client\_backend}, \texttt{src/data\_api/chain.rs:598--699}).
This is
precisely the deployed cost centre that the \emph{Tachyon Guide}'s
proof-carrying wallet is designed to remove (\S``The proof-carrying
wallet'', cost-centre paragraph ``Trial decryption of the chain'').

\subsection{The compact stream}\label{sec:sync-scan-stream}

ZIP~307 (Status: Draft) defines the light-client protocol for payment
detection. Its \emph{compact output} retains, of a full shielded output, the
$32$-byte note commitment, the $32$-byte ephemeral key, and the first $52$
bytes of the note ciphertext:
\[
\underbrace{32}_{\text{commitment}} + \underbrace{32}_{\mathsf{epk}} +
\underbrace{52}_{\text{ciphertext prefix}} = 116 \text{ bytes},
\]
an ${\sim}80\%$ reduction against the full $580$-byte ciphertext plus
$32$-byte ephemeral key; its \emph{compact spend} is the $32$-byte nullifier
alone, an ${\sim}90\%$ reduction against the $384$-byte Sapling Spend
description (ZIP~307 \S\S``Output Compression'', ``Spend Compression'').
Truncating at $52$ bytes discards the AEAD tag, so a compact decryption
authenticates differently from a full one: the wallet recomputes the note
commitment from the decrypted compact plaintext and compares it with the
transmitted commitment, trusting the transaction assembler for integrity
(ZIP~307 \S``Compact Stream Format''). The cryptography of trial decryption
--- key agreement, KDF, the $52$-byte compact note plaintext, the commitment
recomputation --- is the \emph{Orchard Guide}'s material (\S``Transmission
and detection of a note: encryption, trial decryption, and
synchronisation''); this volume treats it as a black box mapping a (compact
output, incoming viewing key) pair to a note or $\bot$.

\paragraph{The deployed wire format.} The format in production has evolved
well beyond ZIP~307's message definitions. A block is
\begin{center}
\texttt{CompactBlock \{ reserved 1; height=2; hash=3; prevHash=4; time=5;}\\
\texttt{header=6; vtx=7; chainMetadata=8 \}}
\end{center}
where field $1$ (\texttt{protoVersion}) has been removed and
\texttt{ChainMetadata} carries the running per-pool note commitment tree
sizes \texttt{saplingCommitmentTreeSize=1},
\texttt{orchardCommitmentTreeSize=2}, and
\texttt{ironwoodCommitmentTreeSize=3} --- the field on which position
tracking (\S\ref{sec:sync-scan}) depends. Each \texttt{CompactTx} carries
its index, txid, fee, Sapling \texttt{spends} and \texttt{outputs}, Orchard
\texttt{actions}, Ironwood actions (\texttt{ironwoodActions}, field $9$,
reusing the \texttt{CompactOrchardAction} message), and transparent
\texttt{vin}/\texttt{vout}
(\textsf{zcash\_client\_backend}, \texttt{lightwallet-protocol/}%
\allowbreak\texttt{walletrpc/compact\_formats.proto}, vendored at
\textsf{librustzcash} \texttt{e30517e4}).

\begin{remark}[Normative anchoring]\label{rem:sync-spec-of-record}
ZIP~307's message layout --- a \texttt{CompactBlock} with an embedded
\texttt{BlockID}, messages \texttt{CompactSpend} and \texttt{CompactOutput}
--- does not match the deployed format, and the ZIP remains Draft and
Sapling-only. No ZIP specifies the deployed compact format or service:
\texttt{ChainMetadata} tree sizes, \texttt{CompactOrchardAction},
\texttt{ironwoodActions}, \texttt{GetTreeState}, and
\texttt{GetSubtreeRoots} exist only in the \texttt{zcash/lightwallet-protocol}
proto files, which we therefore cite as the specification of record.
Ironwood's compact-format extensions are likewise absent from ZIP~2005,
which covers quantum recoverability, not light-client synchronisation
(\texttt{compact\_formats.proto}; ZIP~307; ZIP~2005).
\end{remark}

\paragraph{The fee field.} The \texttt{fee} of a \texttt{CompactTx} is
calculable by a stateless server only for transactions without transparent
inputs, as
\[
\begin{aligned}
\mathsf{fee} = {} & \mathsf{valueBalanceSapling}
+ \mathsf{valueBalanceOrchard} + \mathsf{valueBalanceIronwood} \\
& + \textstyle\sum \mathsf{vPubNew} - \sum \mathsf{vPubOld}
- \sum \mathsf{tOut};
\end{aligned}
\]
with transparent inputs the fee depends on prior-transaction outputs a
stateless server cannot provide (\texttt{compact\_formats.proto:63--69}).

\paragraph{Service surface.} The reference synchroniser drives four RPCs of
the \texttt{CompactTxStreamer} service:
\texttt{GetLatestBlock(ChainSpec) $\to$ BlockID};
\texttt{GetBlockRange(BlockRange) $\to$ stream CompactBlock}, addressed by
heights only, with an optional \texttt{poolTypes} pruning filter that
clients MUST verify the server supports;
\texttt{GetTreeState(BlockID) $\to$ TreeState}, returning network, height,
hash, time, and the hex-serialised Sapling, Orchard, and Ironwood
commitment trees; and
\texttt{GetSubtreeRoots(GetSubtreeRootsArg) $\to$ stream SubtreeRoot},
returning complete-subtree roots with their completing block hash and
height, selected by \texttt{enum ShieldedProtocol \{ sapling=0; orchard=1;
ironwood=2 \}}
(\texttt{lightwallet-protocol/walletrpc/service.proto:30--231}).

\begin{remark}[Deployment divergence at NU6.3]\label{rem:sync-divergence}
The deployed \textsf{lightwalletd} tree (commit \texttt{61fee32e},
2026-06-25) contains no Ironwood support: its \texttt{walletrpc} protos and
Go code lack \texttt{ironwoodActions}, \texttt{ironwoodCommitmentTreeSize},
\texttt{ironwoodTree}, and \texttt{ShieldedProtocol.ironwood}. The protos
vendored by \textsf{librustzcash} (commit \texttt{e30517e4}) and
\textsf{Zaino} (commit \texttt{0e057e22}) define all four, and the reference
synchroniser unconditionally requests Ironwood subtree roots
(\textsf{zcash\_client\_backend}, \texttt{src/sync.rs:279--296}). An NU6.3-aware
wallet therefore requires Zaino or an updated lightwalletd.
%% TODO: determine and document the reference client's observed behaviour
%% against a stock lightwalletd: GetSubtreeRoots with
%% shieldedProtocol=ironwood, and proto-default zero
%% ironwoodCommitmentTreeSize after NU6.3 activation (candidate
%% TreeSizeInvalid/TreeSizeUnknown failures, cf. \S\ref{sec:sync-scan}).
\end{remark}

\subsection{The account birthday}\label{sec:sync-birthday}

\begin{definition}[Account birthday]\label{def:sync-birthday}
Type \texttt{AccountBirthday} pairs a \texttt{prior\_chain\_state:
ChainState} with a \texttt{recover\_until: Option<BlockHeight>}. The
\emph{birthday height} is
\[
\mathsf{birthday} := \mathsf{priorChainState.height} + 1,
\]
the height of the first block to be scanned
(\textsf{zcash\_client\_backend}, \texttt{src/data\_api.rs:3013--3106}). A
\texttt{ChainState} holds a block height, a block hash, and the final note
commitment tree frontier of each pool: Sapling (depth
\texttt{sapling::NOTE\_COMMITMENT\_TREE\_DEPTH} $= 32$), Orchard, and
Ironwood --- the Ironwood tree is Orchard-shaped
(\texttt{MerkleHashOrchard}, the same depth) but the pool is distinct and
keeps its own tree
(\textsf{zcash\_client\_backend}, \texttt{src/data\_api/chain.rs:504--596}).
\end{definition}

Constructor \texttt{AccountBirthday::from\_treestate} builds the birthday
from a lightwalletd \texttt{TreeState} for the last block \emph{prior} to
the birthday height. The conversion (\texttt{TreeState::to\_chain\_state})
hex-decodes the block hash --- reversing the bytes, since zcashd renders
block hashes byte-reversed in hex --- parses the height, and converts the
hex-serialised Sapling, Orchard, and Ironwood trees into per-pool
frontiers; an empty tree field yields an empty tree
(\textsf{zcash\_client\_backend}, \texttt{src/data\_api.rs:3055--3073};
\texttt{src/proto.rs:367--463}).

\paragraph{Choosing a birthday.} For a new wallet the crate documentation
directs constructing the birthday from the treestate at
$\mathsf{tip} - 100$: placing the birthday below the pruning depth
guarantees that no reorganisation can mine a transaction intended for the
wallet \emph{below} its birthday and so cause funds to be missed
(\textsf{zcash\_client\_backend}, \texttt{src/data\_api.rs:3199--3211}). For a
restored wallet the birthday is historical, and the \texttt{recover\_until}
height marks the chain tip at restoration time: the wallet is \emph{in
recovery} until no unscanned ranges remain between the birthday height and
\texttt{recover\_until} (exclusive). The field exists to avoid confusing
shifts of balance and spendability while restore-from-seed is in flight ---
spendability policy itself is the \emph{Wallet Guide}'s
\texttt{ConfirmationsPolicy} (\S``Transaction construction'';
\S``Broadcast, expiry, and the transaction lifecycle''), which this volume
does not restate. Progress reporting follows the same split: type
\texttt{Progress} carries two scanned-notes/notes-on-chain ratios, one over
the recovery window (birthday to recovery height) and one over the scan
window (recovery height to tip); either denominator may be zero
(\textsf{zcash\_client\_backend},
\texttt{src/data\_api.rs:3042--3047,788--831}).

\paragraph{The wallet birthday.} The birthday of the \emph{wallet} is the
minimum birthday over its accounts --- literally
\texttt{SELECT MIN(birthday\_height) FROM accounts} in the SQLite backend
--- and adding an account whose birthday lies below the current chain tip
triggers a rescan of all blocks at and above that height
(\textsf{zcash\_client\_sqlite}, \texttt{src/wallet.rs:2967--2978};
\textsf{zcash\_client\_backend}, \texttt{src/data\_api.rs:3153--3156}). Function
\texttt{add\_account} persists the birthday height together with the
Sapling and Orchard frontier tree sizes, calls
\texttt{rewind\_to\_chain\_state} targeting the birthday's prior chain
state --- installing a \texttt{Historic}-priority rescan range above
$\mathsf{birthday} - 1$ up to the wallet's known chain tip while preserving
higher-priority queue entries --- and finally inserts an
\texttt{Ignored}-priority range covering
$[\mathsf{saplingActivation}, \mathsf{birthday})$
(\textsf{zcash\_client\_sqlite}, \texttt{src/wallet.rs:432--628,4056--4224}).

\subsection{The scan queue}\label{sec:sync-queue}

The wallet does not scan the chain linearly; it maintains a priority queue
of height ranges and always works on the most urgent.

\begin{definition}[Scan range]\label{def:sync-scanrange}
A \texttt{ScanRange} is a half-open range $[\mathsf{start}, \mathsf{end})$
of block heights paired with a \texttt{ScanPriority}, with helpers
\texttt{truncate\_start}, \texttt{truncate\_end}, and \texttt{split\_at}
(\textsf{zcash\_client\_backend}, \texttt{src/data\_api/scanning.rs:31--125}).
The
SQLite backend stores the queue in table \texttt{scan\_queue} as rows
\texttt{(block\_range\_start, block\_range\_end, priority)} with both
bounds \texttt{UNIQUE} and a \texttt{CHECK} constraint
$\mathsf{start} < \mathsf{end}$
(\textsf{zcash\_client\_sqlite}, \texttt{src/wallet/db.rs:1017--1028}).
\end{definition}

\begin{table}[htbp]
\centering
\begin{tabular}{@{}llp{8.2cm}@{}}
\toprule
Priority & Code & Meaning (doc comment) \\
\midrule
\texttt{Ignored}      & 0  & Lowest priority; never scanned. \\
\texttt{Scanned}      & 10 & Already scanned; not re-scanned. \\
\texttt{Historic}     & 20 & Advance the fully-scanned height. \\
\texttt{OpenAdjacent} & 30 & Ranges adjacent to heights at which the user
                             opened the wallet (vestigial;
                             Remark~\ref{rem:sync-openadjacent}). \\
\texttt{FoundNote}    & 40 & Complete tree shards adjacent to found notes. \\
\texttt{ChainTip}     & 50 & Complete the latest shard. \\
\texttt{Verify}       & 60 & Previously scanned range that must be
                             re-verified as still on the main chain
                             (highest). \\
\bottomrule
\end{tabular}
\caption{The \texttt{ScanPriority} ladder, ascending, with its SQLite codes
(\textsf{zcash\_client\_backend}, \texttt{src/data\_api/scanning.rs:11--29};
\textsf{zcash\_client\_sqlite}, \texttt{src/wallet/scanning.rs:34--45}).}
\label{tab:sync-priorities}
\end{table}

\begin{remark}[Six live priorities plus one vestigial]
\label{rem:sync-openadjacent}
We present the ladder as six live priorities and one vestige. Priority
\texttt{OpenAdjacent} is declared and mapped to SQL code $30$, but no code
path in the current \textsf{librustzcash} tree ever assigns it: a search at
commit \texttt{e30517e4} finds only the enum declaration, the code mapping,
and two doc-comment mentions
(\textsf{zcash\_client\_backend}, \texttt{src/data\_api/scanning.rs:20--21}). It
still participates in the ordering, so we retain it in
Table~\ref{tab:sync-priorities}.
\end{remark}

\paragraph{Suggesting work.} Method \texttt{WalletRead::suggest\_scan\_ranges}
returns queue entries in descending priority order. Its documented contract
has two clauses: \texttt{Verify} ranges must be scanned before all others,
to avoid continuity errors under reorganisation; and the compact-block
source must supply note commitment tree size metadata
(\textsf{zcash\_client\_backend}, \texttt{src/data\_api.rs:1995--2010}). The
SQLite
implementation orders by \texttt{priority DESC, block\_range\_end DESC} ---
nearest the tip first within a priority --- and filters to priorities at or
above \texttt{Historic}
(\textsf{zcash\_client\_sqlite}, \texttt{src/wallet/scanning.rs:61--90};
\textsf{zcash\_client\_sqlite}, \texttt{src/lib.rs:958--960}).

\begin{remark}[Contract, not type]\label{rem:sync-verify-first}
The guarantee that a \texttt{Verify} range arrives first is a documented
contract of the trait realised by the SQLite \texttt{ORDER BY}; nothing in
the type system enforces it, and a third-party \texttt{WalletRead}
implementation could violate it. We state it as a trait obligation and cite
the one deployed implementation that discharges it
(\textsf{zcash\_client\_backend}, \texttt{src/data\_api.rs:1995--2010};
\textsf{zcash\_client\_sqlite}, \texttt{src/wallet/scanning.rs:61--90}).
\end{remark}

\paragraph{Queue maintenance.} Function \texttt{replace\_queue\_entries}
performs every insertion as a spanning-tree merge: existing entries
overlapping or adjacent to the query range are deleted and re-inserted
together with the new entries, and overlaps are resolved by a dominance
rule ---
\begin{itemize}
\item an inserted \texttt{Verify} or \texttt{Scanned} range always replaces
  what it overlaps;
\item otherwise an existing \texttt{Scanned} range wins (unless a forced
  rescan is requested);
\item otherwise the higher priority wins, and equal priorities coalesce
  into one range.
\end{itemize}
Gaps between disjoint ranges are filled at \texttt{Historic} priority, so
the queue always tiles a contiguous span
(\textsf{zcash\_client\_sqlite}, \texttt{src/wallet/scanning.rs:145--222};
\textsf{zcash\_client\_backend}, \texttt{src/data\_api/scanning/}%
\allowbreak\texttt{spanning\_tree.rs:65--153}).

\subsection{Tracking the chain tip}\label{sec:sync-tip}

Function \texttt{update\_chain\_tip} feeds the queue whenever the wallet
learns a new tip height $\mathsf{tip}$. Two derived heights recur, with
range upper bounds exclusive throughout:
\[
\mathsf{chainEnd} := \mathsf{tip} + 1, \qquad
\mathsf{stable} := \mathsf{tip} - \mathsf{PRUNING\_DEPTH},
\]
where $\mathsf{PRUNING\_DEPTH} = 100$ blocks and heights at or below
$\mathsf{stable}$ are treated as not subject to reorganisation
(\textsf{zcash\_client\_sqlite},
\texttt{src/wallet/scanning.rs:522--523,582--584};
\textsf{zcash\_client\_sqlite}, \texttt{src/lib.rs:169}). The SQLite
implementation
proceeds as follows
(\textsf{zcash\_client\_sqlite}, \texttt{src/wallet/scanning.rs:488--660}):
\begin{enumerate}[label=(\alph*)]
\item it is a no-op if the tip is below Sapling activation;
\item it is a no-op if the new tip is below the prior maximum scanned
  height --- the caller has caught the chain mid-reorganisation, and the
  discontinuity error path of \S\ref{sec:sync-reorg} will resolve it;
\item it computes a \emph{tip shard entry} at \texttt{ChainTip} priority
  covering every pool's last incomplete shard,
  \[
  \Bigl[\,\max\bigl(\mathsf{birthday},\
  \min_{\text{pools}} \max(\mathsf{subtreeEndHeight})\bigr),\
  \mathsf{chainEnd}\Bigr),
  \]
  taking the minimum over the Sapling, Orchard, and Ironwood shards tables
  and emitting the entry only if that minimum lies below
  $\mathsf{chainEnd}$;
\item it computes a \emph{tip entry}: with no scanned blocks, a
  \texttt{Historic} range from the wallet birthday to $\mathsf{chainEnd}$
  (or \texttt{Ignored} from Sapling activation when no accounts exist);
  with no shard metadata (linear scanning), \texttt{Historic} from
  $\mathsf{maxScanned} + 1$; if $\mathsf{maxScanned} > \mathsf{stable}$
  (steady state), \texttt{ChainTip} from $\mathsf{maxScanned} + 1$ to
  $\mathsf{chainEnd}$; otherwise a \texttt{Verify} range
  \[
  \bigl[\,\mathsf{minUnscanned},\
  \min(\mathsf{stable} + 1,\ \mathsf{minUnscanned} +
  \mathsf{VERIFY\_LOOKAHEAD})\,\bigr),
  \]
  where $\mathsf{minUnscanned} := \mathsf{maxScanned} + 1$.
\end{enumerate}

The \texttt{Verify} discipline is documented at the same site: because
\texttt{Verify} outranks everything, the connectivity check --- and any
rewind it forces --- runs before any \texttt{ChainTip} range from this or
an earlier tip update is scanned. Constant
$\mathsf{VERIFY\_LOOKAHEAD} = 10$ blocks is chosen as roughly $12.5$
minutes at $75\,\mathrm{s}$ per block, a window within which a user may
plausibly have received a transaction without leaving the wallet open; and
the \texttt{Verify} range is clamped to the stable region so that it cannot
itself be reorganised away and interfere with later \texttt{Verify} ranges
(\textsf{zcash\_client\_sqlite}, \texttt{src/wallet/scanning.rs:583--634};
\textsf{zcash\_client\_sqlite}, \texttt{src/lib.rs:172}).

\begin{remark}[Two deliberate asymmetries]\label{rem:sync-tip-asymmetries}
First, when $\mathsf{maxScanned} = \mathsf{stable}$ the \texttt{Verify}
range above is empty: after a long offline period no \texttt{Verify} entry
need exist, and connectivity is then checked only implicitly, by the
previous-hash comparison of \S\ref{sec:sync-reorg} when the adjacent range
is scanned (\textsf{zcash\_client\_sqlite},
\texttt{src/wallet/scanning.rs:625--634}).
Second, case (b) above means \texttt{update\_chain\_tip} deliberately does
\emph{nothing} when the tip moves backwards; reorganisation detection lives
in the scanner, not in tip updates, and recovery in that case waits for a
fetch failure or continuity error
(\textsf{zcash\_client\_sqlite}, \texttt{src/wallet/scanning.rs:488--660}).
\end{remark}

\subsection{Scanning a range}\label{sec:sync-scan}

Function \texttt{scan\_cached\_blocks} scans one contiguous range; it
takes the network parameters, the block source, the wallet database, a
starting height \texttt{from\_height} with its prior chain state
\texttt{from\_state}, and an optional block limit. It panics unless
$\mathsf{fromHeight} = \mathsf{fromState.height} + 1$ --- the caller must
supply the chain state immediately prior to the range
(\textsf{zcash\_client\_backend}, \texttt{src/data\_api/chain.rs:598--699}).

\paragraph{Keys.} The scanner loads all tracked unified full viewing keys
and builds a \texttt{ScanningKeys} set
(\texttt{ScanningKeys::from\_account\_ufvks}). Per account it registers:
the Sapling diversifiable full viewing key's (External, Internal) pairs of
incoming viewing key and nullifier-deriving key; the Orchard full viewing
key's (External, Internal) incoming viewing keys; and the \emph{same}
Orchard keys a second time under the Ironwood note-encryption domain ---
Ironwood notes are Orchard-shaped and decrypt under Orchard viewing keys,
but carry version-3 note plaintexts, so \texttt{IronwoodDomain} is distinct
from \texttt{OrchardDomain}, which accepts only version-2 plaintexts. Each
incoming viewing key is tagged \texttt{(AccountId, zip32::Scope)}
(\textsf{zcash\_client\_backend}, \texttt{src/scanning.rs:213--256,352--434};
the
Ironwood pool is ZIP~2005, ``Ironwood Quantum Recoverability'', Status:
Proposed, NU6.3). A scanning
key derives a note's nullifier only when it carries nullifier-deriving
capability: trait method \texttt{ScanningKeyOps::nf(note, position)}
returns \texttt{None} for bare incoming viewing keys, so notes detected
with an IVK alone never enter the tracked-nullifier set and their spends
cannot be detected
(\textsf{zcash\_client\_backend}, \texttt{src/scanning.rs:34--66,168--186}).

\paragraph{Batched trial decryption.} The scanner makes two passes over the
downloaded blocks. The first pass feeds every compact output into a set of
\texttt{BatchRunners} --- one per note-encryption domain: Sapling, Orchard,
Ironwood --- and flushes. A runner accumulates (domain, output) pairs
across transactions and blocks into a \texttt{Batch} shared by all prepared
incoming viewing keys of its domain; when the accumulated output count
reaches the batch-size threshold ($100$ outputs in
\texttt{scan\_cached\_blocks}) the batch is spawned on the global
\textsf{rayon} thread pool (\texttt{rayon::spawn\_fifo}) and a fresh batch
begins. Results are delivered per (block hash, txid) over \textsf{flume}
channels, and \texttt{collect\_results} blocks until the relevant batch has
run (\textsf{zcash\_client\_backend}, \texttt{src/scan.rs:194--220,461--626};
\texttt{src/scanning/compact.rs:101--238};
\texttt{src/data\_api/chain.rs:642}). Within a batch,
\texttt{zcash\_note\_encryption::batch::try\_compact\_note\_decryption}
batches the cryptography: \texttt{D::batch\_epk} parses and prepares all
ephemeral keys at once, shared secrets are computed for every
(\ivk, \textsf{epk}) pair --- the scalar multiplications themselves cannot
be batched --- and \texttt{D::batch\_kdf} derives all symmetric keys in one
pass before per-output trial decryption of the $52$-byte compact ciphertext
(\textsf{zcash\_note\_encryption}, \texttt{src/batch.rs:42--85}); the
underlying primitives are again the \emph{Orchard Guide}'s
(\S``Transmission and detection of a note: encryption, trial decryption,
and synchronisation'').

\paragraph{The second pass.} The scanner then reads the prior block
metadata and the wallet's unspent-nullifier set
(\texttt{Nullifiers::unspent}, per pool), and calls
\texttt{scan\_block\_with\_runners} per block, accumulating summary counts,
updating the in-memory nullifier set after each block
(\texttt{Nullifiers::update\_with}: nullifiers spent in the block are
removed, nullifiers of newly received full-viewing-key-decrypted notes are
added), and chaining each block's metadata into the next --- so a spend of
a note received earlier in the same batch is detected before anything is
persisted (\textsf{zcash\_client\_backend},
\texttt{src/data\_api/chain.rs:598--699};
\texttt{src/scanning.rs:436--600}).

Within a block, spend detection (\texttt{find\_spent}) compares each
revealed nullifier against the tracked unspent set in constant time --- a
\textsf{subtle} \texttt{CtOption} fold over the whole set; an in-code TODO
notes the resulting
$O(|\text{nullifiers}| \times |\text{spends}|)$ cost and questions whether
constant-time operation is warranted. Nullifiers matching no tracked note
are recorded as \emph{unlinked} in a per-transaction nullifier map, in case
they match a note found later when scanning an earlier range
(\textsf{zcash\_client\_backend}, \texttt{src/scanning.rs:826--870}).
%% TODO: the threat model for the constant-time fold (what timing observer
%% it defends against on a local scan) is undocumented in code and ZIPs.
Note detection (\texttt{find\_received}) assigns each decrypted note its
absolute tree position,
$\mathrm{position}(o) = \mathsf{treeSize}_{\text{tx start}} +
\mathrm{idx}(o)$, and computes a \textsf{shardtree} retention for every
commitment in the block:
\[
\mathrm{retention}(c) =
\begin{cases}
\texttt{Checkpoint\{Marked\}} & \text{last in block, wallet's} \\
\texttt{Checkpoint\{None\}}   & \text{last in block only} \\
\texttt{Marked}               & \text{wallet's only} \\
\texttt{Ephemeral}            & \text{otherwise,}
\end{cases}
\]
and a note is flagged \texttt{is\_change} when the receiving account also
spent notes in the same transaction
(\textsf{zcash\_client\_backend}, \texttt{src/scanning.rs:804--824,872--994}).

\paragraph{Position tracking.} Positions require tree sizes, which the
compact stream must supply. Per pool, the tracker
(\texttt{PositionTracker::for\_compact\_block}) reconstructs the block's
starting tree size: from the prior block's metadata when available; else
from the block's \texttt{ChainMetadata} final size minus the block's output
count (error \texttt{TreeSizeInvalid} on underflow --- the proto-default
zero is caught here); else $0$ below the pool's activation height (Sapling,
NU5, and NU6.3 for Ironwood respectively); else \texttt{TreeSizeUnknown}.
The tracker advances by each transaction's per-pool output count, and at
end of block its computed size is checked against \texttt{ChainMetadata}
--- a mismatch is the continuity error \texttt{TreeSizeMismatch}
(\textsf{zcash\_client\_backend}, \texttt{src/scanning/compact.rs:577--791}).
Malformed input from the untrusted server is likewise caught up front:
wrong-length nullifier bytes and malformed compact outputs or actions yield
\texttt{ScanError::EncodingInvalid} --- naming pool, txid, and index ---
rather than a panic
(\textsf{zcash\_client\_backend},
\texttt{src/scanning/compact.rs:174--234,312--390}).

\paragraph{Output.} Each block yields a \texttt{ScannedBlock}: height,
hash, block time, the wallet-relevant \texttt{WalletTx} list --- a
transaction is retained only if it has at least one wallet spend or output
in any pool --- and per-pool \texttt{ScannedBundles} holding the final tree
size, the (commitment, retention) list, and the nullifier map. Transparent
data in a \texttt{CompactTx} is not yet scanned (\textsf{librustzcash}
issue \#2187)
(\textsf{zcash\_client\_backend}, \texttt{src/scanning/compact.rs:240--575};
\texttt{src/data\_api.rs:2520--2692}). The per-call \texttt{ScanSummary}
reports the scanned range and counts of spent and received Sapling and
Orchard notes; received counts may include notes already spent in later
blocks (a scanning-order effect), and spent counts cover only previously
detected notes
(\textsf{zcash\_client\_backend},
\texttt{src/data\_api/chain.rs:437--502,677--685}).

\begin{remark}[A gap in \texttt{ScanSummary}]\label{rem:sync-summary-gap}
No Ironwood counters exist in \texttt{ScanSummary}, although
\texttt{ScannedBlock} and \texttt{WalletTx} carry Ironwood spends and
outputs; Ironwood activity is silently excluded from the summary. We read
this as an oversight in the deployed code rather than a design decision
(\textsf{zcash\_client\_backend},
\texttt{src/data\_api/chain.rs:437--502,677--685};
\texttt{research/upstream-defects.md}, item~2: ``ScanSummary lacks
spent/received counters for the Ironwood pool'', observed at
\textsf{librustzcash} \texttt{e30517e433}).
\end{remark}

The Ironwood pool otherwise traverses the pipeline as a full peer of
Sapling and Orchard: its actions are trial-decrypted under
\texttt{IronwoodDomain} with the account's Orchard viewing keys, its
nullifiers form a separate set --- method \texttt{get\_ironwood\_nullifiers}
is a required trait method, deliberately without a panicking default
because it sits on the scan path --- and its commitment tree and scan-queue
shard extensions are its own
(\textsf{zcash\_client\_backend}, \texttt{src/scanning.rs:213--256};
\texttt{src/data\_api.rs:2066--2078}).

\subsection{Persistence}\label{sec:sync-persist}

The scanned range is committed by a single call,
\texttt{WalletWrite::put\_blocks(from\_state, blocks)}, which requires
sequential blocks in increasing height order with \texttt{from\_state} the
chain state prior to the first block. The backend implementation validates,
for the first block $b_0$ and every pool,
\[
\mathsf{fromState.height} + 1 = \mathsf{height}(b_0)
\quad\text{and}\quad
\mathsf{fromState.treeSize} + |\mathsf{commitments}(b_0)| =
\mathsf{finalTreeSize}(b_0),
\]
each subsequent height being the predecessor plus one; violation is
\texttt{NonSequentialBlocks}. In \textsf{zcash\_client\_sqlite} the whole
call runs in one database transaction
(\textsf{zcash\_client\_backend}, \texttt{src/data\_api.rs:3478--3490};
\texttt{src/data\_api/ll/wallet.rs:291--348};
\textsf{zcash\_client\_sqlite}, \texttt{src/lib.rs:1425--1431}).

\paragraph{Stage 1: rows.} Per block, \texttt{put\_blocks} persists block
metadata (heights, hash, time, per-pool final tree sizes and commitment
counts); per transaction it persists transaction metadata, queues the txid
for full-transaction retrieval (\emph{enhancement}), marks spent notes, and
persists received notes --- checking each received note's nullifier against
the \emph{persisted} nullifier map (\texttt{detect\_*\_spend}) to catch
notes spent in a later block range that was scanned earlier, the
out-of-order counterpart of the in-memory check of \S\ref{sec:sync-scan}.
It then persists each block's unlinked-nullifier map and prunes the map
beyond $\mathsf{PRUNING\_DEPTH}$
(\textsf{zcash\_client\_backend},
\texttt{src/data\_api/ll/wallet.rs:291--579,55}).

\paragraph{Stage 2: trees.} Note commitments are assembled into shard
subtrees in parallel --- chunks of $1024$ commitments, built at the pool's
shard height, which is $16$ for all three pools (subtrees of $2^{16}$
commitments; the Ironwood tree is Orchard-shaped) --- and each shard tree
is updated. The same checkpoint set is enforced across all three trees:
each tree gains a checkpoint at every height checkpointed in any other.
Checkpoints at or above the NU6.3 activation height that fall on the
$\mathsf{ANCHOR\_RETENTION\_INTERVAL} = 288$-block grid are retained as
durable anchors exempt from checkpoint pruning
(\textsf{zcash\_client\_backend},
\texttt{src/data\_api/ll/wallet.rs:597--720,1537};
\texttt{src/data\_api.rs:159,166,173}). This is the boundary of the present
volume: shard-tree structure, witness derivation, and anchor selection are
the \emph{Orchard Guide}'s material (\S``Proof of ownership of \emph{some}
valid note: the commitment tree'').

\paragraph{Completing the found notes' shards.} After the tree update,
\texttt{scan\_complete} marks the scanned range \texttt{Scanned} in the
queue and, if wallet notes were found, extends the queue with
\texttt{FoundNote}-priority ranges covering the blocks needed to complete
each found note's subtree --- the blocks that make the note witnessable.
Per pool, the shard of a note at position $p$ is
\[
\mathsf{subtreeIndex}(p) = \lfloor p / 2^{16} \rfloor
\qquad (\texttt{Address::above\_position(SHARD\_HEIGHT, } p\texttt{)}),
\]
and the extension spans from the previous shard's end height --- bounded
below by the wallet birthday and the pool's activation height --- to the
containing shard's end height plus one. The same pass marks notes
\texttt{witness\_stabilized} once their shard's block extent is fully
\texttt{Scanned} and the shard's end height is at most the pruning floor,
\[
\mathsf{floor} := \mathsf{maxScanned} - (\mathsf{PRUNING\_DEPTH} - 1)
\]
(\textsf{zcash\_client\_sqlite}, \texttt{src/wallet/scanning.rs:224--473}).

\paragraph{Two notions of scanned height.} Out-of-order scanning leaves
gaps, so the backend distinguishes \texttt{block\_max\_scanned} --- simply
$\max(\texttt{blocks.height})$ --- from \texttt{block\_fully\_scanned}, the
height to which this and all preceding blocks above the wallet birthday
have been trial-decrypted. The SQLite implementation computes the latter
from the queue: it is $\mathsf{end} - 1$ of the \emph{first}
\texttt{Scanned} range (by ascending start), defined only when that range
starts at or below the birthday
(\textsf{zcash\_client\_backend}, \texttt{src/data\_api.rs:1974--1993};
\textsf{zcash\_client\_sqlite}, \texttt{src/wallet.rs:3221--3307}).

\subsection{Reorganisation detection and recovery}\label{sec:sync-reorg}

ZIP~307 prescribes the primitive check: on each sync, re-fetch the block at
the previous sync height and compare its hash with the cached copy; a
mismatch means the block was orphaned (ZIP~307 \S``Client-server
interaction'', Phases B and D). The
deployed scanner generalises this into a per-block continuity check.
Function \texttt{scan\_block\_with\_runners} first runs
\texttt{check\_hash\_continuity}, returning
\texttt{ScanError::BlockHeightDiscontinuity} if
$\mathsf{height}(b) \ne \mathsf{height}(\mathrm{prev}) + 1$ and
\texttt{ScanError::PrevHashMismatch} if
$\mathsf{prevHash}(b) \ne \mathsf{hash}(\mathrm{prev})$; for the first
block of a \texttt{scan\_cached\_blocks} call the prior metadata is read
from the wallet database (\texttt{block\_metadata(from\_height - 1)}), and
for subsequent blocks it is the just-scanned block's own metadata
(\textsf{zcash\_client\_backend}, \texttt{src/scanning/compact.rs:257--289};
\texttt{src/data\_api/chain.rs:649--693}). Of ZIP~307's block-header
validation list only this previous-hash check is implemented, and the
ZIP's ``block privacy via bucketing'' enhancement --- rounding the sync
start down to a bucket multiple --- is explicitly unimplemented
(ZIP~307 \S\S``Block header validation'', ``Block privacy via
bucketing'').

Classification is explicit: method \texttt{is\_continuity\_error} on
\texttt{ScanError} returns true for \texttt{PrevHashMismatch},
\texttt{BlockHeightDiscontinuity}, and \texttt{TreeSizeMismatch}; the
encoding and tree-size-availability errors of \S\ref{sec:sync-scan}
(\texttt{EncodingInvalid}, \texttt{TreeSizeUnknown},
\texttt{TreeSizeInvalid}, \texttt{TreeSizeOverflow}) are not continuity
errors (\textsf{zcash\_client\_backend}, \texttt{src/scanning.rs:602--685}).

\paragraph{Recovery.} On a continuity error the reference loop rewinds:
\[
\mathsf{rewind} := \mathsf{errHeight} - 10 \quad
(\texttt{saturating\_sub}),
\]
where any height at least one block before the error is valid and the
$10$-block offset is a heuristic. It then calls
\texttt{truncate\_to\_height(rewind)} on the wallet database,
\texttt{truncate(rewind)} on the block cache, and returns control so the
caller restarts from \texttt{update\_chain\_tip} and a fresh
\texttt{suggest\_scan\_ranges}
(\textsf{zcash\_client\_backend}, \texttt{src/sync.rs:423--453}).
Truncation semantics --- resolution to a shared shard-tree checkpoint at
or below the requested height (error \texttt{RequestedRewindInvalid},
carrying a safe rewind height, otherwise), un-mining of transactions
above it, and the deliberate retention of note rows --- are the
\emph{Wallet Guide}'s material (\S``Reorganisations: truncate and
rescan''). This volume adds two consequences within its own scope: the
scan queue is trimmed to the truncation height, rolling back the
wallet's view of the chain tip so that the next
\texttt{update\_chain\_tip} re-installs \texttt{Verify} and
\texttt{ChainTip} ranges above the new maximum scanned height; and at
the pinned commit truncation covers all three shard trees and the
nullifier-map locators
(\textsf{zcash\_client\_sqlite}, \texttt{src/wallet.rs:3707--3897,4252--4274}).

\paragraph{The depth constant.} Constant $\mathsf{PRUNING\_DEPTH} = 100$
blocks is \textsf{librustzcash}'s reorganisation-safety depth
(\textsf{zcash\_client\_sqlite}, \texttt{src/lib.rs:169}), the nullifier-map
pruning
depth (\textsf{zcash\_client\_backend}, \texttt{src/data\_api/ll/wallet.rs:55}),
and
ZIP~307's suggested witness-cache size --- the ZIP asserts $100$ blocks is
reasonable ``as no full Zcash node will roll back the chain by more than
100 blocks'' (ZIP~307 \S``Client operation''). The claim is grounded in
\textsf{zcashd}: $\mathsf{MAX\_REORG\_LENGTH} =
\mathsf{COINBASE\_MATURITY} - 1 = 99$ (\texttt{src/main.h:66}), and a
node asked to reorganise deeper aborts
(\texttt{src/main.cpp:4332--4337}). \textsf{Zebra}'s rollback window is
$\mathsf{MAX\_BLOCK\_REORG\_HEIGHT} = 1000$, documented as local-only
node policy, not consensus
(\texttt{zebra-chain/src/parameters/constants.rs:20--30}), so the ZIP's
``no full Zcash node'' claim no longer holds for Zebra-backed
deployments; $\mathsf{PRUNING\_DEPTH} = 100$ matches the zcashd bound,
as its doc comment says (\textsf{zcash\_client\_sqlite},
\texttt{src/lib.rs:166--169}).

\subsection{The reference sync loop}\label{sec:sync-loop}

Module \texttt{sync} of \textsf{zcash\_client\_backend} assembles the
pieces into the reference flow, whose contract the chain module documents
(\texttt{src/sync.rs:58--226}, \texttt{393--467};
\texttt{src/data\_api/}\allowbreak\texttt{chain.rs:1--152}):
\begin{enumerate}[label=(\arabic*)]
\item[(1--2)] \texttt{update\_subtree\_roots}: download and store complete
  subtree roots for Sapling, Orchard, and Ironwood via
  \texttt{GetSubtreeRoots} (cf.\ Remark~\ref{rem:sync-divergence});
\item[(3--4)] \texttt{update\_chain\_tip} with the height from
  \texttt{GetLatestBlock};
\item[(5)] \texttt{suggest\_scan\_ranges};
\item[(6)] scan every \texttt{Verify} range first: download it via
  \texttt{GetBlockRange}, fetch the prior \texttt{ChainState} via
  \texttt{GetTreeState} at $\mathsf{start} - 1$, scan, delete the cached
  blocks;
\item[(7)] scan the remaining ranges, split into chunks of the
  caller-supplied batch size, restarting from step (3) whenever scanning
  materially changed the suggestions --- a new first range outranking the
  range just scanned, or a truncation after a continuity error
  (\S\ref{sec:sync-reorg}).
\end{enumerate}

Blocks pass through a \texttt{BlockCache} (an extension of
\texttt{BlockSource}): \texttt{get\_tip\_height}, \texttt{read} --- short
reads allowed but contiguous from the range start --- \texttt{insert} ---
non-contiguity permitted --- \texttt{delete} and \texttt{truncate}. The
loop deletes each range from the cache after scanning it, the module
observing that ``keeping the entire chain in \texttt{CompactBlock} files on
disk is horrendous for the filesystem''; deletions are spawned concurrently
and awaited before the loop returns
(\textsf{zcash\_client\_backend}, \texttt{src/data\_api/chain.rs:237--435};
\texttt{src/sync.rs:131--226}).

The module documents its own limitations: block batches are not downloaded
in parallel with scanning; detected transactions are not enhanced (the full
transaction, hence the memo, is not fetched); there are no progress
notifications; and there is no interruption mechanism short of process
exit (\textsf{zcash\_client\_backend}, \texttt{src/sync.rs:1--10}). The
chain-module
documentation additionally recommends scanning \texttt{Historic} ranges in
reverse height order --- or from both ends inward --- so that recent
unspent notes are discovered sooner
(\textsf{zcash\_client\_backend}, \texttt{src/data\_api/chain.rs:128--134}).

\section{Commitment-tree synchronisation}\label{sec:sync-tree}
%% Ground truth: librustzcash @ e30517e4 (2026-07-10); lightwalletd @
%% 61fee32e (2026-06-25); incrementalmerkletree + shardtree @ edf24f2
%% (2026-02-20); zaino @ 0e057e22 (2026-07-04); zcashd @ 16ac74376
%% (v6.10.0, 2025-10-03); zebra @ 9f3166b96 (2026-07-10); zips @
%% 69610984 (2026-07-09). Line references are against these commits.

Scanning tells a wallet which notes it owns (\S\ref{sec:sync-scanning});
spending one additionally requires the note's authentication path against
a recent anchor. The construction of that path --- the depth-$32$ tree cut
at half depth into height-$16$ shards under a height-$16$ cap, the wallet
hashing its own shard from scanned leaves and taking every other shard as
a single downloaded $32$-byte root --- is the \emph{Orchard Guide}'s
material (\S``Proof of ownership of \emph{some} valid note: the
commitment tree'', paragraph ``Derivation of the authentication path'').
This section supplies what that construction presumes: the RPCs that
deliver tree states and shard roots, their derivation inside the
validator, and the client discipline --- frontier-seeded scanning,
reference-retained roots, shard-completion scan ranges --- by which
\textsf{zcash\_client\_backend} and \textsf{zcash\_client\_sqlite} turn a
compact-block stream plus two RPCs into spendable witnesses.

\subsection{The linear protocol and its limit}\label{sec:sync-tree-linear}

ZIP~307 (Status: Draft) specifies only the original, linear protocol: as
compact blocks arrive in height order, the client appends every note
commitment to a local depth-$32$ incremental Merkle tree and updates one
incremental witness per unspent note; because incremental trees cannot be
rewound, the ZIP advises caching roughly $100$ recent tree and witness
states, on the stated ground that no full node rolls back the chain by
more than $100$ blocks (\texttt{zips/zip-0307.rst:290--320}; the same
figure resurfaces as $\mathsf{PRUNING\_DEPTH}$, \S\ref{sec:sync-reorg}).
Under this design a wallet cannot spend a note until it has scanned every
block from the note's block to the anchor block --- each witness update
consumes every intervening commitment --- so restore-from-seed implies a
complete linear scan before the first spend. Removing that
spend-before-sync impossibility is the purpose of the machinery below.

\begin{remark}[Normative status]\label{rem:sync-tree-normative}
No ZIP specifies \texttt{GetTreeState}, \texttt{GetSubtreeRoots}, or the
subtree-root chain index. ZIP~307's Phase~A leaves tree-state acquisition
an explicit unresolved TODO --- ``Or, it has to set X to the block height
at which Sapling activated, so as to be sent the entire commitment tree.
[TODO: Decide which to specify.]''
(\texttt{zips/zip-0307.rst:433--439}) --- and the repository holding the
canonical protobufs disclaims normativity: ``This repository IS not a
specification of the Zcash Light Client protocol''
(\texttt{lightwallet-protocol/README.md}, v0.4.1). As in
Remark~\ref{rem:sync-spec-of-record}, we therefore cite the proto files
and the deployed implementations as the specification of record, and
classify the subtree-root protocol \emph{deployed but unspecified}.
\end{remark}

\subsection{The tree-state service surface}\label{sec:sync-tree-rpcs}

\begin{definition}[Tree-state and subtree-root RPCs]\label{def:sync-tree-rpcs}
The deployed light-wallet service defines three tree RPCs
(\texttt{lightwallet-protocol/walletrpc/}\allowbreak
\texttt{service.proto:307--316}):
\begin{itemize}
\item \texttt{GetTreeState(BlockID) returns (TreeState)} --- the tree
  state at a block specified by height or by hash;
\item \texttt{GetLatestTreeState(Empty) returns (TreeState)} --- the
  tree state at the server's chain tip;
\item \texttt{GetSubtreeRoots(GetSubtreeRootsArg) returns (stream
  SubtreeRoot)} --- a stream of information about roots of subtrees of
  the note commitment tree of the requested shielded protocol ---
  Sapling or Orchard in the deployed protocol.
\end{itemize}%
\footnote{The proto comment refers the reader to ``section 3.7 of the
Zcash protocol specification''
(\texttt{service.proto:308}). The reference is stale: at the pinned spec
revision, \S3.4 is ``Transactions and Treestates'', \S3.8 ``Note
Commitment Trees'', and \S3.7 ``Action Transfers and their
Descriptions''
(\texttt{protocol/protocol.tex:4048,\allowbreak 4248,\allowbreak 4306}),
likely a
pointer into an older revision's numbering. We cite the current
numbering.}
\end{definition}

\begin{definition}[\texttt{TreeState} message]\label{def:sync-treestate}
Message \texttt{TreeState}, documented in the proto as ``derived from
the Zcash \texttt{z\_gettreestate} rpc'', carries
\[
\begin{aligned}
&\mathsf{network} : \text{string} = 1
  \ (\texttt{"main"} \mid \texttt{"test"}); \quad
 \mathsf{height} : \text{uint64} = 2; \\
&\mathsf{hash} : \text{string} = 3
  \ (\text{big-endian display-order hex}); \quad
 \mathsf{time} : \text{uint32} = 4; \\
&\mathsf{saplingTree} : \text{string} = 5; \quad
 \mathsf{orchardTree} : \text{string} = 6,
\end{aligned}
\]
where the two tree fields are protobuf \emph{strings} --- hex encodings
of the legacy-serialised commitment trees of
Definition~\ref{def:sync-legacy-ser} --- not \texttt{bytes}
(\texttt{lightwallet-protocol/walletrpc/}\allowbreak
\texttt{service.proto:176--184}). At
\textsf{librustzcash} HEAD the vendored proto adds
$\mathsf{ironwoodTree} : \text{string} = 7$, absent from the deployed
v0.4.1 surface (Remark~\ref{rem:sync-divergence}).
\end{definition}

\begin{definition}[Legacy commitment-tree encoding]\label{def:sync-legacy-ser}
The tree state of one pool at one block is serialised in the legacy
\texttt{CommitmentTree} wire format,
\[
\mathrm{ser}(T) \;=\; \mathrm{Opt}(\mathsf{left}) \,\|\,
\mathrm{Opt}(\mathsf{right}) \,\|\,
\mathrm{Vec}\bigl(\mathrm{Opt}(\mathsf{parents}_1), \mathrm{Opt}(
\mathsf{parents}_2), \ldots\bigr),
\]
where $\mathrm{Opt}(\cdot)$ writes $\mathtt{0x00}$ for an absent node and
$\mathtt{0x01}$ followed by the $32$-byte node otherwise,
$\mathrm{Vec}(\cdot)$ is CompactSize-prefixed, and each node is the
$32$-byte encoding of a base-field element --- Jubjub base for Sapling,
Pallas base for Orchard --- with non-canonical encodings rejected at
parse time. The parser \texttt{read\_commitment\_tree} bounds the
$\mathsf{parents}$ vector at $\mathrm{DEPTH} - 1$ entries, and the parsed
tree converts losslessly to a frontier via \texttt{to\_frontier()}
(\textsf{zcash\_primitives}, \texttt{src/merkle\_tree.rs:26--61,183--209}).
The producer is \textsf{zcashd}: the Sapling final state streams the
legacy \texttt{SaplingMerkleTree} directly and the Orchard final state
streams through the \texttt{OrchardMerkleFrontierLegacySer} wrapper, so
both pools use the same wire format, hex-encoded
(\textsf{zcashd},
\texttt{src/rpc/blockchain.cpp:1402--1406,\allowbreak 1432--1436}).
\end{definition}

\begin{definition}[Subtree-root messages]\label{def:sync-subtreeroot}
Message \texttt{GetSubtreeRootsArg} carries
$\mathsf{startIndex} : \text{uint32} = 1$, the index of the first subtree
to return; $\mathsf{shieldedProtocol} = 2$ with
\texttt{enum ShieldedProtocol \{ sapling = 0; orchard = 1 \}} as
deployed; and $\mathsf{maxEntries} : \text{uint32} = 3$, where $0$ means
``return all entries''. Results stream in index --- equivalently height
--- order. Message \texttt{SubtreeRoot} carries
$\mathsf{rootHash} : \text{bytes} = 2$, the $32$-byte Merkle root of the
completed subtree; $\mathsf{completingBlockHash} : \text{bytes} = 3$; and
$\mathsf{completingBlockHeight} : \text{uint64} = 4$. Field number $1$ is
absent, and no \texttt{reserved} statement documents the gap
(\texttt{lightwallet-protocol/walletrpc/}\allowbreak
\texttt{service.proto:186--200}).
\end{definition}

\begin{remark}[Byte-order warts]\label{rem:sync-tree-byteorder}
Two byte-order pitfalls attach to these messages. First,
\texttt{SubtreeRoot.completingBlockHash} is sent in big-endian display
order --- \textsf{lightwalletd} sets it to
\texttt{hash32.ReverseSlice(block.Hash)} --- which is the \emph{opposite}
convention from \texttt{CompactBlock.hash}, parsed into little-endian
wire order; the proto documents neither convention
(\textsf{lightwalletd}, \texttt{frontend/service.go:896--900};
\texttt{parser/block.go:59--61,115}). The reference client sidesteps the
trap by ignoring \texttt{completingBlockHash} entirely, reading only
\texttt{rootHash} and \texttt{completingBlockHeight}
(\textsf{zcash\_client\_backend}, \texttt{src/sync.rs:247--253}).
Zaino constructs the field as the block hash's
\texttt{bytes\_in\_display\_order()}
(\texttt{zaino-state/src/}\allowbreak\texttt{indexer.rs:854--861}),
agreeing with
\textsf{lightwalletd}'s big-endian display order; whether any deployed
client consumes \texttt{completingBlockHash} remains an open question
(the reference client ignores it).
Second, a \textsf{zcashd} release note records that the serialised
\texttt{finalRoot} field of \texttt{z\_gettreestate} was byte-flipped
relative to \texttt{getrawtransaction} output in releases up to and
including v5.2.0, later rectified
(\textsf{zcashd}, \texttt{src/rpc/blockchain.cpp:1326--1330}, help text).
\end{remark}

\subsection{Server-side derivation}\label{sec:sync-tree-server}

\paragraph{Tree states.} The upstream RPC is \textsf{zcashd}'s
\texttt{z\_gettreestate}: it accepts a block hash or height (negative
heights allowed, $-1$ denoting the last known valid block), requires the
block to be on the main chain, and returns the block's hash, height, and
time together with one object per pool,
$\{\mathsf{skipHash}, \mathsf{commitments} : \{\mathsf{finalRoot},
\mathsf{finalState}\}\}$. The $\mathsf{finalState}$ --- the serialised
tree of Definition~\ref{def:sync-legacy-ser} --- is present only when the
block's anchor is retrievable from the coins view
(\texttt{GetSaplingAnchorAt} / \texttt{GetOrchardAnchorAt}); otherwise
$\mathsf{skipHash}$ names the most recent prior block that has one,
stopping at the pool's activation height
(\textsf{zcashd}, \texttt{src/rpc/blockchain.cpp:1303--1455}).

The proxy wraps this in a retry loop. Handler \texttt{GetTreeState} of
\textsf{lightwalletd} rejects a request specifying neither height nor
hash, marshals a height as a decimal string and a hash as big-endian
hex, and calls \texttt{z\_gettreestate}; if the reply's Sapling
$\mathsf{finalState}$ is empty and its Sapling $\mathsf{skipHash}$ is
nonempty, it re-issues the RPC at the $\mathsf{skipHash}$, iterating
until a final state is found or no skip hash remains, and failing with
``\texttt{z\_gettreestate did not return treestate}'' otherwise
(\textsf{lightwalletd}, \texttt{frontend/service.go:311--386}). The walk
terminates in practice because \textsf{zcashd} hard-stops it at the
Sapling activation height (the \texttt{saplingActive} lambda;
\textsf{zcashd}, \texttt{src/rpc/blockchain.cpp:1408--1418}).
Handler \texttt{GetLatestTreeState} calls \texttt{getblockchaininfo},
takes its
\texttt{Blocks} value as the latest height, and delegates to
\texttt{GetTreeState} at that height
(\textsf{lightwalletd}, \texttt{frontend/service.go:388--401}).

\begin{remark}[The skip walk keys on Sapling only]\label{rem:sync-skiphash}
Only the Sapling $\mathsf{skipHash}$ is consulted: when the walk lands on
an earlier block, the returned \texttt{TreeState} carries that earlier
block's height, hash, \emph{and Orchard tree} --- not the Orchard state
of the requested block
(\textsf{lightwalletd}, \texttt{frontend/service.go:311--386}).
%% TODO: characterise when the walk is reachable on a healthy mainnet
%% node (pre-Sapling heights error out at the zcashd stop; are there
%% post-NU5 blocks whose Sapling anchor is absent from the coins view?)
%% before presenting the skip semantics as reachable rather than
%% vestigial.
\end{remark}

\paragraph{Subtree roots.} The upstream RPC is
\texttt{z\_getsubtreesbyindex "pool" start\_index (limit)}, added in
\textsf{zcashd} v5.6.0 and gated behind
\texttt{-experimentalfeatures} $+$ \texttt{-lightwalletd}: it returns
roots of complete $2^{16}$-leaf subtrees of the given pool's note
commitment tree, each with $\mathsf{end\_height}$, the height of the
block containing the note commitment that completed the subtree, reading
\texttt{SubtreeData} records from the coins view
(\texttt{pcoinsTip->GetSubtreeData(pool, index)}) until data is absent or
the limit is reached
(\textsf{zcashd}, \texttt{src/rpc/blockchain.cpp:1457--1539};
\texttt{doc/release-notes/}\allowbreak\texttt{release-notes-5.6.0.md:28}).
The records are
maintained by the validator's incremental tree: \textsf{zcashd} stores
$\texttt{SubtreeData}\{\mathsf{leadbyte}, \mathsf{root},
\mathsf{nHeight}\}$ and
$\texttt{LatestSubtree}\{\mathsf{leadbyte}, \mathsf{index},
\mathsf{root}, \mathsf{nHeight}\}$ records under a \texttt{uint64\_t}
subtree index, fixes $\mathsf{TRACKED\_SUBTREE\_HEIGHT} = 16$, and
produces a completed subtree root exactly when the tree size crosses a
$2^{16}$ boundary (\textsf{zcashd},
\texttt{src/zcash/}\allowbreak
\texttt{IncrementalMerkleTree.hpp:20--70,\allowbreak 409--411};
\texttt{IncrementalMerkleTree.cpp:913--925}):
\[
\texttt{current\_subtree\_index()} \;=\; \mathsf{size}() \gg
\mathsf{TRACKED\_SUBTREE\_HEIGHT}.
\]

Handler \texttt{GetSubtreeRoots} of \textsf{lightwalletd} accepts only
Sapling or Orchard and forwards to
\texttt{z\_getsubtreesbyindex(pool, startIndex[, maxEntries when > 0])};
for each returned subtree it fetches the block at
$\mathsf{end\_height}$, populates the two completing-block fields from
it, and streams the assembled \texttt{SubtreeRoot}
(\textsf{lightwalletd}, \texttt{frontend/service.go:832--907}).

\begin{example}[Mainnet Sapling subtrees $0$ and $1$]
A worked datum recorded in \textsf{lightwalletd} source (a quoted
\texttt{z\_getsubtreesbyindex} transcript,
\texttt{common/common.go:178--210}): Sapling subtree $0$ --- outputs $0$
through $65535$ --- has root
\[
\texttt{754bb593ea42d231a7ddf367640f09bbf59dc00f2c1d2003cc340e0c016b5b13}
\]
with $\mathsf{end\_height} = 558822$, and subtree $1$ has root
\[
\texttt{03654c3eacbb9b93e122cf6d77b606eae29610f4f38a477985368197fd68e02d}
\]
with $\mathsf{end\_height} = 670209$: measured from Sapling activation
at height $419{,}200$ (\textsf{zcash\_protocol},
\texttt{src/consensus.rs:488}), the first $2^{17}$ Sapling outputs
span roughly the chain's first $250{,}000$ post-activation blocks.
\end{example}

\paragraph{Alternative back ends.} Zebra implements both
\texttt{z\_gettreestate} and \texttt{z\_getsubtreesbyindex} in its RPC
surface, so either validator can back an indexer for tree
synchronisation (\textsf{zebra-rpc},
\texttt{src/methods.rs:361,2018}; \texttt{src/methods/trees.rs}). The
successor indexer serves the same gRPC API: Zaino's
\texttt{get\_tree\_state} proxies \texttt{z\_get\_treestate} to the
backing validator and hex-encodes the final states (at HEAD it also
populates an \texttt{ironwood\_tree} field), its
\texttt{get\_latest\_tree\_state} reads the chain height and delegates,
and its \texttt{get\_subtree\_roots} fetches each completing block via
\texttt{z\_get\_block(end\_height)}
(\textsf{zaino-state},
\texttt{src/backends/}\allowbreak\texttt{fetch.rs:1836--1896};
\texttt{src/indexer.rs:756--860};
\texttt{src/chain\_index.rs:509--514,\allowbreak 2397--2431}). One behavioural
divergence: Zaino maps the proto's \texttt{uint32}
$\mathsf{startIndex}$/$\mathsf{maxEntries}$ down to \texttt{u16} and
returns \texttt{InvalidArgument} when either exceeds $65535$ ---
semantically safe, since a depth-$32$ tree has at most $2^{16}$ shards,
but stricter than \textsf{lightwalletd}/\textsf{zcashd}, whose
\texttt{u64} subtree index yields an empty stream instead
(\texttt{src/chain\_index.rs:509--514} versus
\textsf{zcashd} \texttt{src/zcash/}\allowbreak
\texttt{IncrementalMerkleTree.hpp:20--70}).
Both backends delegate to the validator:
\texttt{NodeBackedChainIndex::}\allowbreak\texttt{get\_subtree\_roots}
calls the blockchain
source (\texttt{chain\_index.rs:2421--2431});
\texttt{ValidatorConnector::State} issues
\texttt{ReadRequest::}\allowbreak\texttt{\{Sapling,}\allowbreak%
\texttt{Orchard,}\allowbreak\texttt{Ironwood\}}\allowbreak%
\texttt{Subtrees} to Zebra's
\texttt{ReadStateService}, and \texttt{::Fetch} proxies
\texttt{z\_getsubtreesbyindex}
(\texttt{chain\_index/source/}\allowbreak
\texttt{validator\_connector.rs:545--608}); only
the mockchain test source self-computes
(\texttt{chain\_index/source/}\allowbreak
\texttt{mockchain\_source.rs:579--665}).

\subsection{Client state: shards, cap, and retention}
\label{sec:sync-tree-store}

The client-side container is \textsf{shardtree}: ``a sparse binary
Merkle tree of the specified depth, represented as an ordered collection
of subtrees (shards) of a given maximum height'', plus a \emph{cap} ---
the tree above the shard roots --- and a checkpoint set enabling
truncation. The root sits at address $(\mathrm{DEPTH}, 0)$, shards at
level $\mathrm{SHARD\_HEIGHT}$, and
$\texttt{max\_subtree\_index()} =
2^{\mathrm{DEPTH} - \mathrm{SHARD\_HEIGHT}} - 1$
(\textsf{shardtree}, \texttt{src/lib.rs:61--129}). The wallet API
instantiates one such tree per pool as
$\texttt{ShardTree}\langle \mathsf{Store}, \mathrm{DEPTH} = 32,
\mathrm{SHARD\_HEIGHT} = 16 \rangle$ with block heights as checkpoint
identifiers, via \texttt{with\_sapling\_tree\_mut} /
\texttt{with\_orchard\_tree\_mut}
(\textsf{zcash\_client\_backend}, \texttt{src/data\_api.rs:3743--3826}).
The shard height is documented as chosen to conform ``to the structure
of subtree data returned by lightwalletd when using the GetSubtreeRoots
GRPC call'': $\mathsf{SAPLING\_SHARD\_HEIGHT} =
\mathsf{ORCHARD\_SHARD\_HEIGHT} = 16$, half the tree depth of each pool
(\texttt{src/data\_api.rs:155--166}). The geometry --- height-$16$
shards of $2^{16}$ leaves under the $16$-level cap, at most $2^{16}$
shards --- is the \emph{Orchard Guide}'s construction (\S``Proof of
ownership of \emph{some} valid note: the commitment tree''); the
deployed constants realise it exactly (\textsf{shardtree},
\texttt{src/lib.rs:111--129};
\textsf{zcash\_client\_sqlite}, \texttt{src/wallet/db.rs:1446--1447}).
The SQLite backend constructs each tree as
$\texttt{ShardTree::new}(\mathsf{store}, \mathsf{PRUNING\_DEPTH})$ with
$\mathsf{PRUNING\_DEPTH} = 100$: at most $100$ checkpoints --- one per
scanned block boundary --- are retained before pruning, $100$ blocks
being the assumed maximum reorganisation depth and ZIP~307's cache-size
guidance (\S\ref{sec:sync-tree-linear})
(\textsf{zcash\_client\_sqlite},
\texttt{src/lib.rs:169,\allowbreak 2500,\allowbreak 2526,\allowbreak
2558}).

What survives pruning is governed by per-node \emph{retention}
(\textsf{incrementalmerkletree}, \texttt{src/lib.rs:76--107}):
\begin{itemize}
\item \texttt{Ephemeral} --- prunable together with its siblings;
\item \texttt{Checkpoint\{id, marking\}} --- the position is retained
  (its value prunable) while the checkpoint lives, decaying to
  \texttt{Marked}, \texttt{Reference}, or \texttt{Ephemeral} per the
  marking when the checkpoint is removed;
\item \texttt{Marked} --- retained until explicitly removed; the wallet's
  own note positions;
\item \texttt{Reference} --- retained until overwritten by a node with
  \texttt{Ephemeral}, \texttt{Checkpoint}, or \texttt{Marked} retention;
  downloaded subtree roots and seed frontiers.
\end{itemize}

\subsection{Ingesting subtree roots}\label{sec:sync-tree-roots}

The reference sync driver begins \emph{every} run by refreshing the
shard roots: function \texttt{update\_subtree\_roots} downloads all
subtree roots --- the default \texttt{GetSubtreeRootsArg}, i.e.\
$\mathsf{startIndex} = 0$, $\mathsf{maxEntries} = 0$ --- for Sapling and,
with the \texttt{orchard} feature, Orchard, converting each streamed
message by
\[
\texttt{CommitmentTreeRoot::from\_parts}\bigl(
\mathsf{completingBlockHeight},\ H\texttt{::read}(\mathsf{rootHash})
\bigr)
\]
and passing the batches to \texttt{put\_sapling\_subtree\_roots} and
\texttt{put\_orchard\_subtree\_roots}, each with start index $0$
(\textsf{zcash\_client\_backend},
\texttt{src/sync.rs:80--86,\allowbreak 228--299}).
Type \texttt{CommitmentTreeRoot} pairs the subtree end height with the
root hash (\texttt{src/data\_api/}\allowbreak\texttt{chain.rs:180--212}),
and the trait
documentation (\texttt{src/data\_api.rs:3743--3826}) fixes the meaning:
each value ``should be the Merkle root of a subtree \ldots\ containing
$2^{\mathrm{SHARD\_HEIGHT}}$ note commitments''.

The SQLite backend implements both puts via one generic procedure,
\texttt{put\_shard\_roots} parameterised by the hash $H$, the tree
depth, and the shard height
(\textsf{zcash\_client\_sqlite},
\texttt{src/wallet/}\allowbreak\texttt{commitment\_tree.rs:1107--1227}):
\begin{enumerate}[label=(\arabic*)]
\item no-op on empty input;
\item load the cap and treat it as a standalone tree of
  $\mathrm{DEPTH} - \mathrm{SHARD\_HEIGHT}$ levels whose level-$0$ leaves
  are shard roots: the \texttt{LevelShifter} wrapper offsets the hash,
  \[
  \begin{aligned}
  \mathsf{emptyLeaf} &= H.\mathsf{emptyRoot}(\mathrm{SHARD\_HEIGHT}), \\
  \mathsf{combine}(\ell, a, b) &=
    H.\mathsf{combine}(\ell + \mathrm{SHARD\_HEIGHT}, a, b),
  \end{aligned}
  \]
  and the downloaded roots are \texttt{batch\_insert}ed at
  $\texttt{Position::from}(\mathsf{startIndex})$ with
  \texttt{Retention::Reference}
  (\texttt{src/wallet/}\allowbreak\texttt{commitment\_tree.rs:1122--1181});
\item write the cap back;
\item check contiguity over
  $[\mathsf{startIndex}, \mathsf{startIndex} + n)$:
  \texttt{check\_shard\_discontinuity} rejects a call whose shard-index
  range is neither overlapping nor immediately adjacent to the stored
  $[\min(\mathsf{shardIndex}), \max(\mathsf{shardIndex}) + 1)$ range,
  raising \texttt{Error::SubtreeDiscontinuity} --- subtree roots must
  arrive without gaps
  (\texttt{src/wallet/}\allowbreak\texttt{commitment\_tree.rs:481--520});
\item upsert each root into
  \texttt{\{pool\}\_tree\_shards(shard\_index, subtree\_end\_height,
  root\_hash, shard\_data)}; on conflict only the end height and root
  hash are updated, so existing scanned shard data is never clobbered,
  and the fresh \texttt{shard\_data} written for new rows is a single
  root leaf with \texttt{RetentionFlags::EPHEMERAL}.
\end{enumerate}
The downloaded root is cached beside any persisted subtree rather than
merged eagerly --- ``we want to avoid deserializing the subtree just to
annotate its root node, so we simply cache the downloaded root alongside
of any already-persisted subtree'' --- and the read path resolves the
pair: \texttt{get\_shard} deserialises \texttt{shard\_data} and applies
$\texttt{reannotate\_root}(\mathsf{Some}(\mathsf{rootHash}))$
(\texttt{src/wallet/}\allowbreak
\texttt{commitment\_tree.rs:1190--1193,\allowbreak 412--445}).

\begin{remark}[Idempotent full re-download]\label{rem:sync-tree-refetch}
The driver fixes $\mathsf{startIndex} = 0$ on every run although
\texttt{put\_shard\_roots} supports arbitrary start indices; the upsert
path makes the repetition idempotent, at the cost of re-transferring one
$32$-byte root (plus completing-block metadata) per completed shard per
pool on every sync
(\textsf{zcash\_client\_backend},
\texttt{src/sync.rs:80--86,\allowbreak 228--299}).
As of 2026-07-15 (mainnet tip $3{,}413{,}101$): $1{,}128$ complete
Sapling subtrees (last completing at block $3{,}390{,}009$) plus $765$
Orchard ($3{,}388{,}160$), so each sync re-downloads $1{,}893$ roots ---
$60{,}576$ bytes of root data, $136{,}378$ bytes of \texttt{SubtreeRoot}
messages, $145{,}843$ bytes with the five-byte gRPC frame per stream
element (computed by script; counts cross-check against
\texttt{ChainMetadata} at block $3{,}413{,}000$:
$\lfloor 73{,}932{,}454/2^{16}\rfloor = 1{,}128$,
$\lfloor 50{,}191{,}264/2^{16}\rfloor = 765$). No source documents the
fixed $\mathsf{startIndex}=0$ as intentional.
\end{remark}

\subsection{Frontiers as scan seeds}\label{sec:sync-tree-frontier}

A \texttt{TreeState} becomes client state through
\texttt{TreeState::to\_chain\_state()}: it hex-decodes the block hash and
reverses its bytes (``Zcashd hex strings for block hashes are
byte-reversed''), and parses each pool's legacy tree
(Definition~\ref{def:sync-legacy-ser}) into a frontier; an \emph{empty}
tree-state string parses as the empty tree, the correct state before a
pool's activation
(\textsf{zcash\_client\_backend}, \texttt{src/proto.rs:367--463}). The
result is
\[
\begin{aligned}
\texttt{ChainState}\{\,&\mathsf{blockHeight};\ \mathsf{blockHash}; \\
&\mathsf{finalSaplingTree} : \texttt{Frontier}\langle 32 \rangle;\
\mathsf{finalOrchardTree} : \texttt{Frontier}\langle 32 \rangle\,\},
\end{aligned}
\]
``the final note commitment tree state for each shielded pool, as of a
particular block height''
(\texttt{src/data\_api/chain.rs:504--596}; a third frontier exists at
HEAD, Remark~\ref{rem:sync-divergence}).

Every scanned batch is seeded with the chain state of the block
immediately preceding it, and the seed is verified against the server's
own metadata before any tree is touched:
\[
\mathsf{fromHeight} = \mathsf{seed}.\mathsf{blockHeight} + 1
\quad\text{(asserted)}, \qquad
\mathsf{seed}.\mathsf{treeSize}_P + |\mathsf{cm}_P(b_0)| =
\mathsf{finalTreeSize}_P(b_0)
\]
per pool $P$, where $b_0$ is the batch's first block and the final
tree size is read from the block's \texttt{chainMetadata}; violation
yields error \texttt{NonSequentialBlocks}
(\textsf{zcash\_client\_backend},
\texttt{src/data\_api/chain.rs:616--635};
\texttt{src/data\_api/ll/wallet.rs:291--348}). The metadata sizes are
what give the scanner absolute leaf positions at all --- without them it
cannot even derive nullifiers (\texttt{TreeSizeUnknown}, a hard error),
and a size contradicted by block contents is \texttt{TreeSizeMismatch},
a continuity error triggering the rewind of \S\ref{sec:sync-reorg}
(\texttt{lightwallet-protocol/walletrpc/}\allowbreak
\texttt{compact\_formats.proto:14--17,\allowbreak 31--40};
\textsf{zcash\_client\_backend},
\texttt{src/scanning.rs:602--685}).

The batch's note commitments are then assembled into located subtrees in
parallel --- \textsf{rayon}, chunks of $1024$ commitments,
\texttt{LocatedTree::from\_iter} at the shard level --- starting at
$\texttt{Position::from}(\mathsf{seed}.\mathsf{treeSize}_P)$, and each
pool's tree is updated by \texttt{update\_tree}
(\textsf{zcash\_client\_backend},
\texttt{src/data\_api/ll/}\allowbreak\texttt{wallet.rs:597--771}):
\begin{enumerate}[label=(\arabic*)]
\item insert the seed frontier (\texttt{tree.insert\_frontier}) with
  retention
  $\texttt{Checkpoint}\{\mathsf{id} = \mathsf{frontierHeight},\
  \mathsf{marking} = \texttt{Reference}\}$ --- the in-source comment:
  ``we insert the frontier with Checkpoint retention because we need to
  be able to truncate the tree back to this point'';
\item \texttt{insert\_tree(subtree, checkpoints)} for each built
  subtree;
\item add any cross-pool checkpoints missing above the minimum retained
  checkpoint
  (\texttt{src/data\_api/ll/wallet.rs:1591--1664}; the cross-pool
  checkpoint-unification rule itself, \S\ref{sec:sync-persist}).
\end{enumerate}
Inside \textsf{shardtree}, \texttt{insert\_frontier} splits a non-empty
frontier at the shard boundary: the leaf and the ommers below the shard
level enter the shard containing the frontier's position, the ommers at
levels $\ge \mathrm{SHARD\_HEIGHT}$ enter the cap, checkpoint retention
adds a checkpoint at the leaf position, and checkpoints beyond
\texttt{max\_checkpoints} are pruned. An empty frontier with checkpoint
retention records a checkpoint for the empty tree
(\texttt{Retention::Marked} is invalid there)
(\textsf{shardtree}, \texttt{src/lib.rs:323--405}).

Frontier seeding is also how a wallet is born and how it rewinds. An
\texttt{AccountBirthday} wraps ``the chain state prior to the birthday
height'' plus an optional \texttt{recover\_until};
\texttt{AccountBirthday::from\_treestate} builds it directly from a
\texttt{TreeState} for the last block \emph{prior} to the birthday, with
$\texttt{height()} = \mathsf{priorChainState}.\mathsf{blockHeight} + 1$,
and account creation stores the birthday tree sizes and rewinds the
wallet to the birthday chain state, so no tree information before the
birthday is ever needed
(\textsf{zcash\_client\_backend}, \texttt{src/data\_api.rs:3010--3106};
\textsf{zcash\_client\_sqlite}, \texttt{src/wallet.rs:438--618};
\S\ref{sec:sync-birthday}). Symmetrically,
\texttt{truncate\_to\_chain\_state} inserts each pool's frontier from
the target chain state with
$\texttt{Checkpoint}\{\mathsf{id} = \mathsf{targetHeight},
\mathsf{marking} = \texttt{None}\}$, creating the checkpoint on which the
subsequent tree truncation lands
(\textsf{zcash\_client\_sqlite}, \texttt{src/wallet.rs:3984--4034};
recovery flow, \S\ref{sec:sync-reorg}).

\subsection{Shard completion and spendability}
\label{sec:sync-tree-witness}

Two of the scan-queue priorities of \S\ref{sec:sync-queue} exist purely
for this layer: \texttt{ChainTip} is documented as ``blocks that must be
scanned to complete the latest note commitment tree shard'' and
\texttt{FoundNote} as ``blocks that must be scanned to complete note
commitment tree shards adjacent to found notes''
(\textsf{zcash\_client\_backend},
\texttt{src/data\_api/scanning.rs:13--29}).

\paragraph{The tip shard.} On each tip update the backend computes per
pool the maximum \texttt{subtree\_end\_height} over the pool's shards
table --- the end of the last complete shard known from
\texttt{GetSubtreeRoots} --- takes the minimum across pools, and
enqueues
\[
\bigl[\,\max(\mathsf{walletBirthday},\ \mathsf{minShardTip}),\
\mathsf{tip} + 1\bigr)
\]
at \texttt{ChainTip} priority, so the incomplete tip shard of every pool
is scanned even before historic ranges
(\textsf{zcash\_client\_sqlite},
\texttt{src/wallet/}\allowbreak\texttt{scanning.rs:475--657}; the full
tip-update case
analysis, \S\ref{sec:sync-tip}).

\paragraph{Found notes.} After a range is scanned,
\texttt{scan\_complete} maps each detected note position to its shard,
$s = \lfloor p / 2^{16} \rfloor$
(\texttt{Address::above\_position(SHARD\_HEIGHT, position)}), and
extends the scanned range to that shard's block extent --- from the
previous shard's \texttt{subtree\_end\_height} (the pool activation
height for shard $0$, clamped to the wallet birthday) to the containing
shard's \texttt{subtree\_end\_height} $+\,1$ --- enqueuing the
extensions beyond the just-scanned range at \texttt{FoundNote} priority.
Completing the shard's block range is exactly what makes the note's
witness computable
(\textsf{zcash\_client\_sqlite},
\texttt{src/wallet/scanning.rs:224--398}).

\paragraph{Witness derivation.} The query
\texttt{witness\_at\_checkpoint\_id(position, id)} requires
$\mathsf{position} \le \mathsf{checkpoint.position}()$ (else
\texttt{QueryError::NotContained}) and folds, from the note's shard
upward to the root, the root of each sibling address truncated at leaf
count $\mathsf{position}_{\mathrm{as\,of}} + 1$ --- each sibling root
drawn from the note's own shard (scanned leaves), another shard's cached
root (\texttt{GetSubtreeRoots},
Definition~\ref{def:sync-subtreeroot}), the cap, or a per-level
empty-root constant beyond the tree size. A pruned node that cannot be
replaced by an empty-root constant yields
\texttt{QueryError::}\allowbreak\texttt{TreeIncomplete(addresses)} ---
the precise failure
mode of an unready witness
(\textsf{shardtree}, \texttt{src/lib.rs:1236--1362};
\texttt{src/error.rs:124--136}). The mathematics of the fold is the
\emph{Orchard Guide}'s (\S``Proof of ownership of \emph{some} valid
note: the commitment tree''); what this volume adds is the deployed
readiness condition under which it cannot fail.

\paragraph{Spendability.} That condition is enforced in SQL. Writing
$U_P$ for the rows of view
\texttt{v\_\{pool\}\_shard\_unscanned\_ranges}, note selection excludes
any note at position $p$ for which some
$(\mathsf{s}, \mathsf{e}, \mathsf{h}_0, \mathsf{h}_1) \in U_P$ satisfies
\[
\mathsf{s} \le p < \mathsf{e}
\ \wedge\
\mathsf{h}_0 \le \mathsf{anchorHeight}
\ \wedge\
\mathsf{h}_1 > \mathsf{walletBirthday},
\]
i.e.\ the note's shard has an unscanned block range at or below the
anchor; the same view drives \texttt{unscanned\_tip\_exists}, which
invalidates anchors whose containing shard is not fully scanned
(\textsf{zcash\_client\_sqlite},
\texttt{src/wallet/common.rs:142--162,\allowbreak 762--822}). The views compute
$\mathsf{startPosition} = \mathsf{shardIndex} \ll 16$ and
$\mathsf{endPositionExclusive} = (\mathsf{shardIndex} + 1) \ll 16$, set
the shard's block extent to begin at the previous shard's end height
(the pool activation height for the first shard), join the scan-queue
rows whose block ranges overlap that extent, and filter to priority
above \texttt{Scanned}
(\textsf{zcash\_client\_sqlite}, \texttt{src/wallet/db.rs:1435--1494}).
Unwinding the predicate recovers the informal statement: a note is
spendable at a given anchor exactly when its own shard is scanned up to
the anchor, every earlier shard is covered by a downloaded root, and the
blocks from the last complete shard boundary to the anchor have been
scanned (the \texttt{ChainTip} range above). Which anchor a transaction
then selects, and at what confirmation depth, is the \emph{Wallet
Guide}'s material (\S``The proposal''; \S``Confirmations, trust, and
spendability'').

At the pinned \textsf{librustzcash} commit --- part of the NU6.3 work of
Remark~\ref{rem:sync-divergence}, not of any release deployed with
\textsf{lightwalletd} v0.4.1 --- notes additionally gain a
\texttt{witness\_stabilized} flag once their shard's block extent is
fully \texttt{Scanned} and the shard's end height has at least
$\mathsf{PRUNING\_DEPTH}$ confirmations
($\mathsf{floor} = \mathsf{maxScanned} - (\mathsf{PRUNING\_DEPTH} - 1)$);
stabilised notes bypass the confirmations check at spend time, and the
tip shard by definition can never stabilise
(\textsf{zcash\_client\_sqlite},
\texttt{src/wallet/scanning.rs:416--473};
\texttt{src/wallet/common.rs:700--724}).

\paragraph{What this buys, and what remains.} The linear protocol's
impossibility is gone: a freshly restored wallet downloads one frontier
(its birthday tree state) plus one $32$-byte root per completed shard,
and a note found during scanning becomes spendable as soon as its own
shard's block extent and the tip range are scanned --- not after the
entire chain history. What remains is the trial decryption of every
block from the birthday forward (\S\ref{sec:sync-scanning}) and the
shard-completion scanning of this section; both are cost centres the
\emph{Tachyon Guide}'s proof-carrying wallet is designed to remove
(\S``What the fold removes'', paragraphs ``Trial decryption of the
chain'' and ``Per-note witness maintenance'').

\section{What the server learns}\label{sec:sync-privacy}
%% Ground truth: zips @ 69610984 (2026-07-09); librustzcash @ e30517e433
%% (2026-07-10); lightwalletd @ 61fee32 (2026-06-25); zaino @ 0e057e22
%% (2026-07-04). Body only; assumes the orchard-guide.tex preamble.

A light client outsources block storage and filtering to a server it does
not trust with its keys, then reconstructs its own state by trial
decryption on the device (\emph{Orchard Guide}, \S``Transmission and
detection of a note: encryption, trial decryption, and synchronisation'').
The keys never travel; the \emph{requests} do. This section inventories
what the deployed server-side observer learns from those requests: the
stated threat model, the baseline that genuinely holds, the leaks in the
deployed request pattern, the operator's instruments for recording them,
and the mitigations that actually ship. Every claim about deployed
behaviour cites the implementation at the commits pinned above; where the
code diverges from ZIP text, we flag the divergence in place.

\subsection{The stated security model}\label{sec:sync-privacy-model}

ZIP~307 names one primary goal for the light-client protocol.

\begin{definition}[Payment detection privacy]\label{def:sync-pdp}
The serving proxy must not learn which transactions are addressed to a
given light wallet. Under the additional assumption of network privacy
(``via Tor or similar''), the proxy must also be unable to link different
connections or queries as deriving from the same wallet
(ZIP~307 \S``Security Model'').
\end{definition}

The adversary of ZIP~307 is the node/proxy combination, modelled as
\emph{honest but curious}: it is trusted for chain-state correctness ---
the client accepts the block data it serves --- but assumed to record and
analyse everything it sees. The ZIP explicitly does not address spending
notes privately (ZIP~307 \S``Security Model''); transaction submission
nonetheless flows through the same server, and we treat it in
\S\ref{sec:sync-privacy-submit}.

\begin{remark}[The de facto wire specification]\label{rem:sync-zip307-draft}
ZIP~307 has Status Draft --- it was never finalised --- even though the
protocol it sketches is the deployed synchronisation layer. The deployed
wire format has moved past the sketch: the ZIP's \texttt{CompactBlock}
carries \texttt{vtx} at field~3 and a \texttt{BlockID}, while the deployed
message carries \texttt{vtx} at field~7 plus \texttt{protoVersion},
\texttt{prevHash}, \texttt{time}, an (always unset) \texttt{header}, and
\texttt{ChainMetadata}; the deployed \texttt{CompactTx} adds \texttt{txid},
\texttt{fee}, Orchard actions, and transparent \texttt{vin}/\texttt{vout}
counts (ZIP~307 \S``Compact Stream Format'' vs.\ \textsf{lightwalletd},
\texttt{lightwallet-protocol/walletrpc/compact\_formats.proto}). The
\textsf{lightwallet-protocol} proto repository, not the ZIP, is the de
facto wire specification. The planned in-protocol successor is emptier
still: ZIP~314, ``Privacy upgrades to the Zcash light client protocol'',
has Status Reserved --- a title with no content
(\texttt{zips/zip-0314.rst}). The server implementation delegates likewise:
the \textsf{lightwalletd} README directs developers to the external wallet
app threat model for ``important information about the security and privacy
limitations of light wallets that use lightwalletd''; the repository itself
contains no privacy analysis (\textsf{lightwalletd}, \texttt{README.md}).
\end{remark}

\subsection{What never leaves the device}\label{sec:sync-privacy-local}

The baseline holds. No RPC in the deployed service carries key material:
across the full \texttt{CompactTxStreamer} surface (twenty methods), every
request consists only of heights, hashes, txids, transparent addresses, raw
transaction bytes, and pool or subtree selectors (\textsf{lightwalletd},
\texttt{lightwallet-protocol/walletrpc/service.proto}). Incoming viewing
keys remain on the device: trial decryption of compact outputs is performed
client-side by \texttt{scan\_block} over locally cached
\texttt{CompactBlock}s (\textsf{zcash\_client\_backend},
\texttt{src/scanning.rs}), by the mechanism of the \emph{Orchard Guide},
\S``Transmission and detection of a note: encryption, trial decryption, and
synchronisation''; witness maintenance over the resulting positions is that
of \S``Proof of ownership of \emph{some} valid note: the commitment tree''.
Nullifier matching against the \texttt{CompactSpend} and
\texttt{CompactOrchardAction} nullifier fields is likewise local.

The compact form is the reason a server mediates at all: the per-output
and per-spend reductions, and the recomputed-commitment integrity
substitute for the discarded AEAD tag, are those of
\S\ref{sec:sync-compact} (Table~\ref{tab:sync-compact-sizes};
\S\ref{sec:sync-compact-decrypt}).

\begin{remark}[The transparent boundary]\label{rem:sync-transparent}
The baseline above is a shielded-layer property. The same service carries
transparent-address RPCs whose requests contain cleartext addresses ---
\texttt{GetTaddressTxids} and \texttt{GetTaddressTransactions} (an address
plus a block range), \texttt{GetTaddressBalance(Stream)} (an address list),
\texttt{GetAddressUtxos(Stream)} (an address list plus a start height) ---
so for the transparent layer the server learns the wallet's addresses
outright, by declaration rather than inference (\textsf{lightwalletd},
\texttt{walletrpc/service.proto}). We treat that
request family as out of scope except where the shielded sync driver itself
uses it (\S\ref{sec:sync-privacy-schedule}) and where deployed mitigations
attach to it (\S\ref{sec:sync-privacy-mitigations}).
\end{remark}

\subsection{The observable request schedule}\label{sec:sync-privacy-schedule}

What the server sees is a request schedule. ZIP~307 describes it in four
phases: at first start, \textsf{A} = \texttt{GetLatestBlock} plus
\texttt{GetBlock} at the tip; on each subsequent synchronisation,
\textsf{B} = \texttt{GetLatestBlock} followed by
$\texttt{GetBlockRange}(X, Y)$ from the last-synced height $X$;
\textsf{C} = \texttt{GetTransaction}(\texttt{txid}) for each transaction of
interest after detection; and \textsf{D} = the next
\texttt{GetLatestBlock} plus $\texttt{GetBlockRange}(Y, Z)$
(ZIP~307 \S``Client--server interaction'').

The deployed driver refines this but does not change its shape. Per cycle,
function \texttt{run} of \textsf{zcash\_client\_backend} (module
\texttt{sync}, feature \texttt{sync}) issues: \texttt{GetSubtreeRoots} per
shielded pool, \texttt{GetLatestBlock}, \texttt{GetAddressUtxosStream} over
\emph{all} transparent receivers of every account (cleartext,
Remark~\ref{rem:sync-transparent}), \texttt{GetTreeState} for the block
before each scan range, and \texttt{GetBlockRange} for each suggested scan
range in \texttt{batch\_size} chunks (\textsf{zcash\_client\_backend},
\texttt{src/sync.rs}). At the pinned commit the subtree-root fetch covers
Sapling, Orchard, and Ironwood. Ironwood appears in no light-client ZIP
--- ZIP~2005 covers quantum recoverability and ZIP~258 only NU6.3
deployment --- the request exists solely because the
\textsf{librustzcash} proto fork defines
\texttt{ShieldedProtocol.ironwood = 2} and the driver requests it
(\texttt{src/sync.rs:289--296}; cf.\
Remark~\ref{rem:sync-spec-of-record}).
The module's own documentation notes that this simple implementation does
not itself enhance transactions; enhancement is part of the documented
flow (the crate-root diagram's ``Enhance transactions'' step) and is driven
by wallets via \texttt{transaction\_data\_requests}
(\textsf{zcash\_client\_backend}, \texttt{src/lib.rs}, \texttt{src/sync.rs}).

The \emph{order} of range requests is itself data. Scan ranges carry a
priority (Table~\ref{tab:sync-priorities}), and
\texttt{suggest\_scan\_ranges} returns them in descending priority, then
descending end height (\S\ref{sec:sync-queue}), so \textsf{Verify} and
\textsf{ChainTip} ranges near the tip are fetched first and
\textsf{FoundNote} ranges pre-empt \textsf{Historic} backfill
(\textsf{zcash\_client\_backend}, \texttt{src/data\_api/scanning.rs};
\textsf{zcash\_client\_sqlite}, \texttt{src/wallet/scanning.rs}). The
privacy consequence of that pre-emption is
Remark~\ref{rem:sync-foundnote}.

\subsection{Leaks in the deployed pattern}\label{sec:sync-privacy-leaks}

ZIP~307 asserts that compact-block synchronisation itself ``leaks no
information about which transactions the light client is interested in''
because the compact stream is a broadcast medium
(ZIP~307 \S``Transaction privacy''). The same document contradicts the
claim eight sections earlier, for the request pattern actually deployed:
``The above interaction reveals to the server at the start of each
synchronisation phase (B and D) the block height which the light client had
previously synchronised to. This is an information leak under our security
model'' (ZIP~307 \S``Block privacy via bucketing''). The broadcast argument
holds only for a client that fetches everything from everyone; the deployed
client fetches specific ranges at specific times and follows detections
with specific txids. We take the leaks in turn.

\paragraph{The first request reveals the birthday.} On account creation,
the scan queue marks the range from Sapling activation up to the account
birthday as \textsf{Ignored} --- never downloaded --- and queues scanning
from the birthday itself (\textsf{zcash\_client\_sqlite},
\texttt{src/wallet.rs}). The lowest height the client ever requests
therefore approximates the account's creation or recovery height, a
lifetime-stable identifier of the wallet.

\paragraph{Each cycle reveals the prior sync height.} This is the leak
ZIP~307 concedes above. Its proposed mitigation is bucketing: for a bucket
size $N$, phases \textsf{B} and \textsf{D} would request
\[
  \texttt{GetBlockRange}\bigl(X - (X \bmod N),\; Y\bigr)
  \quad\text{instead of}\quad
  \texttt{GetBlockRange}(X, Y),
\]
hiding the exact previously-synced height within an $N$-block bucket
(ZIP~307 \S``Block privacy via bucketing''). The ZIP marks this ``a
proposed enhancement that has not been implemented'', and the deployed code
confirms: \texttt{download\_blocks} issues \texttt{GetBlockRange} with
start and end exactly equal to the suggested scan-range bounds, with no
rounding (\textsf{zcash\_client\_backend}, \texttt{src/sync.rs}).

\paragraph{Detection is followed by a fetch.} During compact-block
scanning, function \texttt{put\_blocks} records, per scanned block, exactly
the subset of transactions relevant to this wallet, and queues a retrieval
request for each; \texttt{transaction\_data\_requests} later emits
\textsf{Enhancement}(\texttt{txid}) or \textsf{GetStatus}(\texttt{txid}),
which the sync driver resolves via the \texttt{GetTransaction} RPC
(\textsf{zcash\_client\_backend}, \texttt{src/data\_api.rs},
\texttt{src/data\_api/ll/wallet.rs}; \textsf{zcash\_client\_sqlite},
\texttt{src/wallet.rs}). The set of txids fetched in full therefore
\emph{equals} the wallet's own transaction set, and each fetch follows
immediately after the scan of the containing range. To the server this is
a detection oracle: phase \textsf{C} traffic names the very transactions
that payment detection privacy (Definition~\ref{def:sync-pdp}) is supposed
to conceal, correlated in time with the range that contains them.

Here the code diverges from the ZIP, in the unusual direction: the ZIP is
stricter than the deployment. ZIP~307's full-transaction mitigation is
normative --- the client ``SHOULD obscure the exact transactions of
interest by downloading numerous uninteresting transactions as well'',
``SHOULD download all transactions in any block from which a single full
transaction is fetched'', and ``MUST convey to the user that fetching full
transactions will reduce their privacy''
(ZIP~307 \S``Transaction privacy''). No decoy, chaff, or cover-traffic code
exists anywhere in \textsf{librustzcash} or \textsf{lightwalletd} (negative
grep over both repositories at the pinned commits), and the deployed RPC
surface offers no bulk full-transaction fetch to implement the second
SHOULD with --- \texttt{GetBlock} returns compact data
(\textsf{lightwalletd}, \texttt{lightwallet-protocol/walletrpc/service.proto}).
%% TODO: the MUST clause is UI-level and unverifiable from the local
%% repositories (wallet front-ends are not local sources); record its
%% status once a wallet-app source is in scope.

\paragraph{Enhancement fetches carry no timing decorrelation.} The deployed
request taxonomy is
\begin{align*}
\texttt{TransactionDataRequest} \in \{\;
  &\textsf{GetStatus}(\texttt{txid}),\;
   \textsf{Enhancement}(\texttt{txid}),\\
  &\textsf{TransactionsInvolvingAddress}\{\ldots,
   \texttt{request\_at}, \ldots\},\\
  &\textsf{GetSpendingTx}(\cdot)\;\}
\end{align*}
(\textsf{zcash\_client\_backend}, \texttt{src/data\_api.rs}; the last two
variants are feature-gated). Only \textsf{TransactionsInvolvingAddress}
carries a \texttt{request\_at} field, documented as the instant chosen so
as to ``decorrelate this request to a light wallet server from other
requests made by the same caller'', ignorable when the supplier of chain
data is ``private, trusted-for-privacy''. The shielded-note variants
\textsf{Enhancement} and \textsf{GetStatus} have no analogous jitter: the
SQLite backend emits them with no \texttt{request\_at}
(\textsf{zcash\_client\_sqlite}, \texttt{src/wallet.rs}). The one deployed
timing defence (\S\ref{sec:sync-privacy-mitigations}) protects transparent
address checks, not the shielded fetch-on-detect path.

\begin{remark}[FoundNote clustering --- an inferred leak]
\label{rem:sync-foundnote}
When a note is found, \texttt{scan\_complete} inserts scan ranges at
priority \textsf{FoundNote} extending the scanned range so as to complete
the note-commitment-tree shard(s) containing the found note, and
\texttt{suggest\_scan\_ranges} orders by priority, so the server observes
out-of-order, high-priority \texttt{GetBlockRange} fetches clustered around
the blocks that contain the wallet's own notes
(\textsf{zcash\_client\_sqlite}, \texttt{src/wallet/scanning.rs};
\textsf{zcash\_client\_backend}, \texttt{src/data\_api/scanning.rs}). No
source documents this as a leak --- neither ZIP~307 nor the code
documentation --- and the statement above is this volume's own inference
from the scan-queue code. It requires adversarial review --- in
particular, independent confirmation that \textsf{FoundNote} ranges surface
to the server as distinct requests rather than being absorbed into
contiguous suggested ranges --- before it may be cited as fact.
%% TODO: adversarial review per CONVENTIONS \S Verification; downgrade or
%% promote this remark accordingly.
\end{remark}

\paragraph{Mempool polling fingerprints the session.} Field
\texttt{exclude\_txid\_suffixes} of the \texttt{GetMempoolTx} request
carries suffixes of txids the client already holds, so a mempool-polling
client additionally reveals which mempool transactions it has already
seen, a session-linkable fingerprint (\textsf{lightwalletd},
\texttt{walletrpc/service.proto}). No source documents this as a leak; we
record it for completeness.

\subsection{The submission channel}\label{sec:sync-privacy-submit}

Spending is out of ZIP~307's stated scope
(\S\ref{sec:sync-privacy-model}), but the deployed submission path runs
through the same server. The \texttt{SendTransaction} handler hex-encodes
the received raw transaction and forwards it to the backing node's
\texttt{sendrawtransaction} JSON-RPC (\textsf{lightwalletd},
\texttt{frontend/service.go}), so the operator holds the complete
transaction \emph{before} network broadcast and can link it to the
submitting IP address, the TLS session, its timing, and --- absent circuit
isolation (\S\ref{sec:sync-privacy-mitigations}) --- to the same session's
sync and enhancement traffic. The backing full node also sees the
transaction and the \texttt{GetTransaction} txid pattern, but without
client IPs: \textsf{lightwalletd} terminates the client connections. The
lifecycle the transaction then follows --- store-then-broadcast, expiry,
confirmation counting under the \texttt{ConfirmationsPolicy} --- is that of
the \emph{Wallet Guide}, \S``Broadcast, expiry, and the transaction
lifecycle''.

\subsection{The operator's instruments}\label{sec:sync-privacy-operator}

What the operator \emph{can} record is a structural property, independent
of any logging switch. The gRPC endpoint is operator-terminated TLS, so
every connection exposes the client IP and its open and close times --- a
Prometheus gauge, \texttt{grpc\_server\_connections\_current}, tracks the
live connection count --- and go-grpc-prometheus interceptors on both the
unary and streaming paths export per-method request counts and latency
histograms over the HTTP endpoint (default \texttt{127.0.0.1:9068})
(\textsf{lightwalletd}, \texttt{cmd/connstats.go}, \texttt{cmd/root.go}).

Logging narrows the distance between can and does. Every RPC handler logs
its full request parameters at Debug level --- \texttt{gRPC
GetBlockRange(\%+v)}, \texttt{gRPC GetTransaction(\%+v)}, \texttt{gRPC
SendTransaction(\%+v)}, \texttt{gRPC GetTaddressBalance(\%+v)},
\texttt{gRPC GetAddressUtxos(\%+v)} --- and full responses at Trace level
(\textsf{lightwalletd}, \texttt{frontend/service.go}). The default level is
Info (\texttt{cmd/root.go}), so a single flag, \texttt{--log-level 5},
records every block range, txid, raw transaction, and transparent address
queried, to the default log file \texttt{./server.log}, without client
consent or notice. A separate unary interceptor, enabled by
\texttt{--grpc-logging-insecure} (default \texttt{false}; the flag name
itself marks the hazard), logs \{peer address, method, duration, error\}
per request, and carries the source comment ``TODO: anonymize the
addresses. cryptopan?'' (\textsf{lightwalletd},
\texttt{common/logging/logging.go},
\texttt{cmd/root.go}).\footnote{The interceptor is wired on the unary path
only, so streaming RPCs --- including \texttt{GetBlockRange}, the bulk of
sync traffic --- never pass through it; the flag logs a skewed subset of
traffic. Whether this is intentional is not determinable from the source.}
The project's own architecture document is candid: ``Logging may capture
identifiable user data. It hasn't received any privacy analysis yet and
makes no attempt at sanitization'' (\textsf{lightwalletd},
\texttt{docs/architecture.md}).

Zaino, the successor indexer, changes none of this structurally: it serves
the full \texttt{CompactTxStreamer} surface; every handler logs the method
name
at Info level --- the literal message is \texttt{[TEST] received call},
apparent leftover test instrumentation in production code --- and, under
the \texttt{prometheus} feature, emits per-method request counts, duration
histograms, and error counters. The metric names are
\texttt{zaino.grpc.requests\_total},
\texttt{zaino.grpc.request\_duration\_seconds}, and
\texttt{zaino.grpc.errors\_total}, labelled by method and code, with no
peer labels; no peer-address logging exists in the dispatch path
(\textsf{zaino-serve}, \texttt{src/rpc/grpc/service.rs},
\texttt{src/lib.rs}). Its README acknowledges the
setting: ``Due to current potential data leaks / security weaknesses
\ldots\ there is a need to use anonymous transport protocols (such as Nym
or Tor) to obfuscate clients' identities from Zcash's indexing servers
(Lightwalletd, Zcashd, Zaino)'', and the serve crate's documentation states
the ingestor ``has been built so that other ingestors may be added that use
different transport protocols (Nym, TOR)'' --- a capability anticipated,
not implemented: no Nym, arti, or onion code exists in the workspace
(\textsf{zaino}, \texttt{README.md}; \textsf{zaino-serve},
\texttt{src/lib.rs}; negative grep over \texttt{packages/}). Zaino is
pre-1.0; all of the above
is cited at commit \texttt{0e057e22} and may move.

\subsection{Deployed mitigations}\label{sec:sync-privacy-mitigations}

Two mitigations ship in the client stack. Neither touches the shielded
fetch-on-detect leak of \S\ref{sec:sync-privacy-leaks}.

\paragraph{Transport: Tor.} Crate \textsf{zcash\_client\_backend} ships an
arti-based Tor client behind feature \texttt{tor}. Method
\texttt{Client::connect\_to\_lightwalletd} tunnels the tonic gRPC channel
through Tor (behind the additional feature
\texttt{lightwalletd-tonic-tls-webpki-roots}), and
\texttt{Client::isolated\_client} returns a handle
whose streams never share circuits with the parent, documented for ``when
you want separate parts of your program to each have a tor::Client handle,
but where you don't want their activities to be linkable to one another
over the Tor network'' (\textsf{zcash\_client\_backend},
\texttt{src/tor.rs}, \texttt{src/tor/grpc.rs}). Circuit isolation is
exactly the tool that would sever the submission channel of
\S\ref{sec:sync-privacy-submit} from sync traffic. We classify this as
\emph{capability-in-librustzcash, deployment-status-unknown}: whether any
deployed wallet routes \texttt{CompactTxStreamer} traffic through it by
default --- and whether \texttt{SendTransaction} uses a circuit isolated
from sync --- is not determinable from the local repositories.
%% TODO: wallet apps (Zashi, Zingo SDKs) are not local sources; verify
%% default Tor routing there before strengthening this classification.

\paragraph{Timing: exponential jitter, transparent addresses only.} The
SQLite backend schedules ephemeral-address checks by sampling the next
check time from an exponential distribution,
\[
  t_{\mathrm{next}} := t_{\mathrm{from}} + \Delta, \qquad
  \Delta \sim \mathrm{Exp}(\lambda), \qquad
  \lambda = 1/T, \qquad
  \E[\Delta] = T,
\]
with the expected interval $T$ set to $(24 \cdot 60 \cdot 60)/n$ seconds
for $n$ tracked ephemeral addresses --- one expected check per address per
day --- over a shuffled address order, ``so that we don't always check them
in the same order'' (\textsf{zcash\_client\_sqlite},
\texttt{src/wallet/transparent.rs},
\texttt{src/wallet/transparent/ephemeral.rs}). The intent is explicit in
the API documentation: only one such query should be performed at a time,
``to avoid linking multiple transparent addresses as belonging to the same
wallet in the view of the light wallet server''
(\textsf{zcash\_client\_backend}, \texttt{src/wallet.rs},
\texttt{TransparentAddressMetadata::next\_check\_time}). No analogous
jitter exists for shielded-note enhancement
(\S\ref{sec:sync-privacy-leaks}).

\subsection{Summary, and the structural reading}
\label{sec:sync-privacy-close}

Table~\ref{tab:sync-privacy-summary} collects the observables.

\begin{table}[htbp]
\centering\small
\begin{tabular}{@{}p{0.28\textwidth}p{0.31\textwidth}p{0.31\textwidth}@{}}
\toprule
Observable & Reveals & Mitigation status\\
\midrule
Client IP, connection timing &
wallet identity; activity windows &
Tor capability in \textsf{zcash\_client\_backend}; wallet deployment
unknown\\
\addlinespace
Lowest \texttt{GetBlockRange} start &
account birthday &
none\\
\addlinespace
Range start, each cycle &
prior sync height &
bucketing proposed (ZIP~307), unimplemented\\
\addlinespace
Out-of-order \textsf{FoundNote} ranges &
blocks holding the wallet's notes &
none (inferred; Remark~\ref{rem:sync-foundnote})\\
\addlinespace
\texttt{GetTransaction} after scan &
the wallet's own txids, timed &
decoys normative (SHOULD), unimplemented\\
\addlinespace
\texttt{exclude\_txid\_suffixes} &
mempool transactions already seen &
none\\
\addlinespace
\texttt{SendTransaction} &
full transaction, pre-broadcast, linked to origin &
circuit-isolation capability, deployment unknown\\
\addlinespace
Transparent-address RPCs &
addresses outright &
exponential jitter and shuffling (ephemeral addresses only)\\
\bottomrule
\end{tabular}
\caption{What the server learns from the deployed light-client protocol,
and the status of each mitigation. Sources: \S\ref{sec:sync-privacy-leaks}
through \S\ref{sec:sync-privacy-mitigations}.}
\label{tab:sync-privacy-summary}
\end{table}

The structural reading is that no row of the table is an implementation
accident. Every leak above --- IP and timing exposure, range start heights,
\textsf{FoundNote} clustering, fetch-on-detect, submission linkage --- is a
property of server-mediated compact-block synchronisation itself: of
detection by trial decryption over a server-supplied stream. A client that
must ask for ranges reveals where it stopped; a client that must fetch what
it detected reveals what it detected. Within the deployed protocol the
upgrade path is a stub (ZIP~314, Reserved,
Remark~\ref{rem:sync-zip307-draft}); the exit on offer is architectural.
The \emph{Tachyon Guide} treats exactly this layer as the deployed
protocol's removable cost centre: under the fold the wallet never scans,
and services are asserted ``fully oblivious to their clients' transaction
details'' (\S``What the fold removes''), while the unresolved analogues of
these same leaks --- the sync-service trust model, and the
timing-correlation linkage the designers call potentially fatal --- carry
forward as that volume's \S``The sync-service trust model'' and
\S``Timing-correlation linkage''. The present section is the baseline
against which those designs ask to be measured.

\end{document}
